%% Springer Nature sn-jnl şablonu — Journal of Combinatorial Optimization.
%% Yazar–yıl kaynakça: sınıf seçeneği sn-basic (sn-basic.bst, paper/ içinde).
%% Derleme (TeX Live, paper/ içinden): xelatex main; bibtex main; xelatex main; xelatex main
%% "pdflatex" seçeneği yalnız sınıfın breakurl yüklemesini kapatır; motor XeLaTeX'tir (fontspec).
\documentclass[pdflatex,sn-basic]{sn-jnl}

\usepackage{fontspec}
\setmainfont{Times New Roman}
\usepackage[turkish,shorthands=off]{babel}
\usepackage{amsmath,amssymb}
\usepackage{graphicx}
\usepackage{xcolor}
\usepackage{booktabs,array,tabularx}
\newcolumntype{Y}{>{\centering\arraybackslash}X}
\usepackage{pgfplots}
\pgfplotsset{compat=1.18}
\usepgfplotslibrary{groupplots}
\pgfplotstableset{search path={.,..,../..}}
\usetikzlibrary{arrows.meta}
\definecolor{iyiRenk}{RGB}{0,84,159}
\definecolor{kotuRenk}{RGB}{178,24,43}
\newcommand{\iyi}[1]{\textcolor{iyiRenk}{\textbf{#1}}}
\newcommand{\kotu}[1]{\textcolor{kotuRenk}{\textit{#1}}}
\usepackage{xurl}
\usepackage{algorithm}
\usepackage{algpseudocode}
\floatname{algorithm}{Algoritma}
\algrenewcommand\algorithmicrequire{\textbf{Girdi:}}
\providecommand{\doi}[1]{\url{https://doi.org/#1}}
\emergencystretch=2em
\def\keywordname{Anahtar Kelimeler}
\makeatletter\def\input@path{{../}{./}}\makeatother

\raggedbottom
\setlength\maxdepth{2pt}
\AtBeginDocument{\hypersetup{bookmarksdepth=2}}

\begin{document}

\title[Vuruş hedeflemede öğrenme gerekli mi?]{Vuruşları Hedeflemek İçin Öğrenmek Gerekir mi? 2-opt/Or-opt
Tabanlı Hızlı İteratif Yerel Aramada Eğitimsiz Kenar Uzunluğu Vekillerinin Ön-Kayıtlı Sınaması}

\author*[1]{\fnm{Ufuk} \sur{ASIL}}\email{ufuk.asil@ostimteknik.edu.tr}
\affil*[1]{\orgname{Ostim Technical University},
  \orgaddress{\street{OSTİM, 100. Yıl Bulvarı, 55/F}, \postcode{06374},
  \city{Ankara}, \country{Türkiye}}}

\abstract{%
İteratif yerel aramada (ILS) double-bridge vuruşunun hangi tur kenarlarını kıracağı, gezgin satıcı
probleminde arama verimini belirler. Bu karar için graf sinir ağı (GNN) gibi öğrenilmiş skorlar
önerilmektedir; ancak maliyetlerine değip değmedikleri en basit eğitimsiz alternatife karşı nadiren
sınanır. Hızlı bir 2-opt/Or-opt ILS'te rastgele vuruşu; eğitimsiz, $O(1)$ maliyetli bir kenar
uzunluğu vekilini; LKH-3 turlarıyla eğitilmiş bir GNN'yi; iki çok katmanlı algılayıcıyı (MLP) ve
koşu içinde öğrenen bir modeli karşılaştırdık. Seçimler geliştirme setlerinde yapıldı, protokol
SHA-256 ve zaman damgasıyla donduruldu. Geliştirmede kullanılmamış 54 problemde (10 tohum, 10 sn)
dört ön-kayıtlı hipotezin dördü de desteklendi: vekil, rastgele vuruşa göre gap'i 0,172 puan
düşürdü ($p=3{,}7\times10^{-7}$) ve GNN ile kazanç hedefli MLP'ye karşı 0,05 puanlık sınırla
kötü-değil testlerini geçti. Önceden kaydedilen keşifsel bloklarda vekil, ağın süresi bütçeden
düşüldüğünde yayımlanmış NeuroLKH modelinden ayrışmadı; süre sayılmadığında NeuroLKH problemlerin
çoğunda küçük bir farkla öndeydi. Eşit vuruş sayısında öğrenilmiş skorlar vekilden kötü değildi ve
vuruş başına en iyi sinyal eğitimsiz kNN sırasıydı: eşit süredeki sıralamayı sinyalin niteliğinden
çok maliyeti belirler. Çevrimiçi modelin zararı işaretli ödülden geliyordu; kabule dayalı krediyle
kayboldu. LKH-3 içinde hedeflemenin ölçülebilir bir etkisi olmadı. Hızlı ILS'te vuruş hedefleme
öğrenilmiş model gerektirmedi; öğrenilmiş arama bileşenleri önce en basit eğitimsiz vekillere
karşı ve eşit maliyetle sınanmalıdır.
}
\keywords{gezgin satıcı problemi, iteratif yerel arama, vuruş hedefleme, eğitimsiz vekil, ön-kayıtlı deney, NeuroLKH}
% kaynak: E584, E585, E589, E599-E602, E1106, E1107, E1128-E1131, E1148-E1152, E1172


\maketitle

\section{Giriş}\label{sec:giris}

Gezgin satıcı problemi (traveling salesman problem, TSP), rotalama, devre yerleşimi ve üretim
planlamasındaki pek çok kararın ortak çekirdeğidir \citep{reinelt_1991_tsplib, applegate_2006_tsp}.
Pratikte büyük örnekler yerel aramayla çözülür. Bu ailenin en yaygın iskeleti iteratif yerel
aramadır (iterated local search, ILS): tur bir yerel optimuma indirilir, ardından bir pertürbasyonla
sarsılır ve yeniden onarılır \citep{martin_1992_largestep, lourenco_2010_ils}. TSP'de bu pertürbasyon
çoğunlukla double-bridge vuruşudur (kick). Vuruşun hangi tur kenarlarını kıracağı, aramanın sınırlı
zamanı nereye harcayacağını belirler. Bu nedenle küçük görünen bu seçim, kısa bütçeli bir çözücünün
kalitesini doğrudan etkiler \citep{johnson_1997_case, applegate_2003_clk}.

\begin{figure}[!htbp]
\centering
% MANSET (hero) -- OTOMATIK URETIM: evidence/scripts/build_v3.py (figures/scripts/manset.py)
% Veri: figures/data/manset_diff_gnn.dat, manset_diff_random.dat, manset_vitrin.csv (vitrin_sec.py, kategori A)
% Kanit: E789 (vs GNN), E790 (vs Random); vitrin: evidence/vitrin.json
% vitrin_sec kazanma orani (gurultu |d|>0.0156 pt): 70.4%
\begin{tikzpicture}
\begin{groupplot}[group style={group size=1 by 2, vertical sep=1.25cm},
  width=\textwidth, height=5.4cm, xmin=0, xmax=55, xtick=\empty,
  xlabel={Holdout2 instances, sorted by $\Delta$ (54)}, title style={font=\small, yshift=-1ex},
  tick label style={font=\scriptsize}, label style={font=\small},
  y tick label style={/pgf/number format/fixed, /pgf/number format/precision=1},
  ymajorgrids, grid style={black!10}]
\nextgroupplot[ymin=-1.4, ymax=0.8,
  ylabel={$\Delta$ gap, Excess $-$ GNN [pt]}, title={(a) Excess vs.\ GNN}]
\addplot[ybar, bar width=0.72, fill=iyiRenk, draw=none] table[x=rank, y=win, unbounded coords=jump] {figures/data/manset_diff_gnn.dat};
\addplot[ybar, bar width=0.72, fill=kotuRenk, draw=none] table[x=rank, y=loss, unbounded coords=jump] {figures/data/manset_diff_gnn.dat};
\draw[black!60] (axis cs:0,0) -- (axis cs:55,0);
\node[anchor=north west, font=\small, fill=white, inner sep=2pt, align=left] at (axis description cs:0.01,0.97) {\textcolor{iyiRenk}{\textbf{Excess lower gap: 41/54}}\quad \textcolor{kotuRenk}{\textit{GNN lower gap: 13/54}}};
\node[circle, fill=black, inner sep=1.1pt] (v0) at (axis cs:1,-1.18350) {};
\node[anchor=west, font=\scriptsize, fill=white, inner sep=1pt, draw=black!30] (l0) at (axis cs:19,-1.0751) {h2c02 (clustered): Excess 0.29\%, GNN 1.47\%};
\draw[black!70, thin, -{Latex[length=1.5mm]}] (l0.west) -- (v0);
\node[circle, fill=black, inner sep=1.1pt] (v1) at (axis cs:3,-0.81961) {};
\node[anchor=west, font=\scriptsize, fill=white, inner sep=1pt, draw=black!30] (l1) at (axis cs:19,-0.7861) {eg7146 (real): Excess 1.53\%, GNN 2.35\%};
\draw[black!70, thin, -{Latex[length=1.5mm]}] (l1.west) -- (v1);
\node[circle, fill=black, inner sep=1.1pt] (v2) at (axis cs:17,-0.12025) {};
\node[anchor=west, font=\scriptsize, fill=white, inner sep=1pt, draw=black!30] (l2) at (axis cs:19,-0.4970) {h2u12 (uniform): Excess 0.46\%, GNN 0.58\%};
\draw[black!70, thin, -{Latex[length=1.5mm]}] (l2.west) -- (v2);
\node[anchor=east, font=\scriptsize, fill=white, inner sep=1pt, draw=black!30, text=kotuRenk] (c0) at (axis cs:51.5,0.2803) {counter-example h2c13 (clustered): Excess 2.16\%, GNN 1.53\%};
\draw[black!70, thin, -{Latex[length=1.5mm]}] (c0.east) -- (axis cs:53.6,0.2803);
\nextgroupplot[ymin=-1.5, ymax=0.6,
  ylabel={$\Delta$ gap, Excess $-$ Random [pt]}, title={(b) Excess vs.\ Random}]
\addplot[ybar, bar width=0.72, fill=iyiRenk, draw=none] table[x=rank, y=win, unbounded coords=jump] {figures/data/manset_diff_random.dat};
\addplot[ybar, bar width=0.72, fill=kotuRenk, draw=none] table[x=rank, y=loss, unbounded coords=jump] {figures/data/manset_diff_random.dat};
\draw[black!60] (axis cs:0,0) -- (axis cs:55,0);
\node[anchor=north west, font=\small, fill=white, inner sep=2pt, align=left] at (axis description cs:0.01,0.97) {\textcolor{iyiRenk}{\textbf{Excess lower gap: 44/54}}\quad \textcolor{kotuRenk}{\textit{Random lower gap: 10/54}}};
\end{groupplot}
\end{tikzpicture}

\caption{54 holdout2 probleminde 10~sn gap farkı, sıralı: (a) fazla-uzunluk$-$GNN, (b) fazla-uzunluk$-$rastgele; mavi: vekil lehine. Vitrinler, GNN'nin kendi çeyrekler arası aralığında kaldığı problemlerden katman başına en büyük avantajla seçildi; ölçüm gürültüsünün altındaki farklar beraberlik sayıldığında vekil problemlerin \%70'inde GNN'den iyidir.} % V013
\label{fig:manset}
\end{figure}

Vuruşu rastgele değil de ``şüpheli'' kenarlara yöneltmek doğal bir fikirdir. Öğrenilmiş TSP
yöntemlerinin başarılı bir kolu da tam olarak bunu yapar. Güçlü bir yerel aramanın içine bir model
yerleştirir ve modelin kenar tahminleriyle aramayı yönlendirir: grafik sinir ağı (graph neural
network, GNN) ile güdümlü yerel arama \citep{hudson_2022_gnngls}, öğrenilmiş aday kümeleri
\citep{xin_2021_neurolkh} ve ısı haritası güdümlü arama \citep{qiu_2022_dimes, sun_2023_difusco}
bunun örnekleridir \citep{karimimamaghan_2022_mlmeta}. Bu çizgiyi izleyerek,
Lin--Kernighan--Helsgaun (LKH) çözücüsünün LKH-3 sürümüyle üretilen turlarla eğitilmiş bir GNN'nin turda kalmayacağını öngördüğü kenarları vuruşla kırmak
kolay bir adımdır \citep{helsgaun_2017_lkh3}.

Bu adımın değerini ölçmek ise sanıldığı kadar kolay değildir. Öğrenilmiş bir rehber çoğunlukla
rastgele seçime ya da kendi varyantlarına karşı sınanır; aynı geometriden türetilebilecek en basit
eğitimsiz vekile karşı sınanması seyrektir. Maksimum bağımsız küme probleminde GNN çıktısının
rastgele değerlerle değiştirilmesinin başarımı pek değiştirmediği gösterilmiştir
\citep{bother_2022_wrong}. Büyük ölçekli TSP'de de uzaklıktan türetilen basit bir ısı haritası,
öğrenilmiş olanlarla rekabet edebilmiştir \citep{xia_2024_rethinking}. İkinci bir sorun, tahmin
doğruluğunun arama faydasıyla eş tutulmasıdır. Bir modelin ``yanlış'' kenarları iyi ayırt etmesi,
onları kırmanın turu kısaltacağı anlamına gelmez. Modelin çalışma süresi de aynı bütçeden
ödenmelidir. Son olarak, yöntem ve ayar seçimi ile doğrulama aynı veride yapıldığında olumlu sonuç
kolayca abartılır \citep{hooker_1995_testing, nosek_2018_prereg}.

Bu çalışmada vuruş hedeflemeyi bu üç tuzaktan arındırılmış bir düzenekte inceledik. Önce saniyede
binlerce vuruş atabilen hızlı bir ILS test ortamı kurduk ve bütün model sürelerini bütçeye yükledik.
Ardından çalışmayı üç aşamada yürüttük. Aşama~1'de yavaş bir iskelette gözlediğimiz anlamlı yararın
ölçüm gürültüsü olduğunu gördük ve bu iddiayı geri çektik. Aşama~2'de, ön-kayıtlı bir protokolle 38
yeni çok büyük ölçekli tümleşik devre (very-large-scale integration, VLSI) probleminde GNN
hedeflemesini rastgele vuruşa ve eğitimsiz bir $k$-en yakın komşu (kNN) sırasına karşı sınadık. % E797
Aşama~3 doğrulayıcı aşamadır. Yöntem ve havuz ayarları yalnız geliştirme setlerinde seçildi. Protokol
SHA-256 özetiyle donduruldu ve özet koşudan önce OpenTimestamps'a gönderildi; Bitcoin blok
976258'e bağlanan bu kanıt, protokolün sonuçlar analiz edilmeden önce var olduğunu üçüncü
kişilerce doğrulanabilir kılar. Test, geliştirmede hiç kullanılmamış 54 problemde (holdout2)
koşuldu. % E584, E585, E586

Sınanan adaylardan biri, güdümlü yerel aramanın (guided local search, GLS) ``uzun kenar şüphelidir''
ilkesini \citep{voudouris_1999_gls} vuruş hedeflemeye taşıyan, eğitimsiz ve kenar başına $O(1)$
maliyetli bir fazla-uzunluk vekilidir. Kenar uzunluğuna dayanan basit bir sinyaldir: uzunluktan uç
noktaların yerel komşu ölçeğini çıkarır. Şekil~\ref{fig:manset} sonucu problem düzeyinde gösterir: her çubuk bir problemdir ve
sıfırın altındaki mavi çubuklar vekilin daha düşük gap verdiği problemlerdir. Vekil 54 problemin
41'inde GNN'den, 44'ünde rastgele vuruştan daha iyi sonuç verir; çubukların en uzunları kümeli ve
gerçek problemlerdedir. % E789, E790
GNN'ye karşı ön-kayıtlı iddia bir kötü-değil iddiasıdır; problemlerin çoğundaki üstünlük keşifsel
bir bulgudur. Aynı ölçütü kendi vekilimize de uyguladık: protokol sonrası keşifsel bir kolda yerel
ölçek düzeltmesi olmadan ham kenar uzunluğu da vekilden ortalamada ayrışmadı. Bu, yararın kenar
uzunluğu bilgisinin kendisinden geldiğini düşündürür (Bölüm~\ref{sec:bulgular-duyarlilik}). % E972

Bu çalışmanın başlıca katkıları şunlardır:
\begin{itemize}
  \item Kenar uzunluğuna dayalı eğitimsiz fazla-uzunluk vekilinin geliştirmede kullanılmamış 54 problemde rastgele vuruşa göre ortalama
  gap'i 0,172 puan düşürdüğünü ($p=3{,}7\times10^{-7}$) ve GNN ile kazanç hedefli çok katmanlı
  algılayıcıya (multilayer perceptron, MLP) 0,05 puanlık sınırla kötü olmadığını ön-kayıtlı
  testlerle gösteriyoruz. % E600, E601, E602, E589
  \item Öğrenilmiş skorlarda kestirim doğruluğunun arama faydasına kendiliğinden dönüşmediğini
  ölçüyoruz: Aşama~2'de GNN yanlış kenarları 58 problemin 57'sinde kNN'den isabetli ayırt ettiği
  hâlde aramada onun gerisinde kaldı; holdout2'de ikisi ayrışmadı, ucuz ve sürekli güncel vekil ise
  ikisinden de iyi sonuç verdi (keşifsel). % E890, E805, E605, E604, E607
  \item Koşu içinde çevrimiçi öğrenmenin aynı doğrusal modelin sabit hâline göre gap'i 0,291 puan
  artırdığını ($p=1{,}6\times10^{-10}$) ve mekanizmasını ölçüyoruz: model kendi vuruş ödülünü
  yükseltirken (12/12 problem) yerel aramanın onardığı nötr vuruşlara yönelir (\%47'den \%88'e). % E599, E941, E944
  \item Geliştirmede seçim, SHA-256 ile dondurma, dış zaman damgası ve 54 bağımsız problemle
  doğrulamadan oluşan bir ön-kayıt protokolünü ve ``öğrenilmiş skoru önce en basit eğitimsiz vekile
  karşı sına'' ilkesini sunuyor, ilkeyi kendi vekilimize de uyguluyoruz: daha basit ham kenar
  uzunluğu keşifsel olarak vekilden ayrışmadı. % E585, E586, E972
\end{itemize}

Uygulayıcı için sonuç sadedir: bu motorda ve 10~sn bütçede vuruşları hedeflemek kazandırır;
sınanan öğrenilmiş skorlar karşısında bunun için model eğitmek gerekmedi. % E588

Makalenin geri kalanı şöyle düzenlenmiştir.
Bölüm~\ref{sec:ilgili-calismalar} ilgili çalışmaları, Bölüm~\ref{sec:yontem} test ortamını ve
hedefleme sinyallerini, Bölüm~\ref{sec:deney-kurulum} üç aşamanın protokollerini verir.
Bölüm~\ref{sec:bulgular} doğrulayıcı sonuçları ve mekanizmaları, Bölüm~\ref{sec:tartisma}
yorumları ve sınırlılıkları sunar; Bölüm~\ref{sec:sonuc} çalışmayı bağlar.

\section{İlgili Çalışmalar}\label{sec:ilgili-calismalar}

\subsection{ILS ve vuruşun yeri}
Double-bridge vuruşu \citet{martin_1992_largestep}'in büyük adımlı Markov zincirleriyle yaygınlaşmış,
\citet{lourenco_2010_ils} ise ILS'yi yerel arama, pertürbasyon ve kabul kuralından oluşan bir çatı
olarak sistematikleştirmiştir. Hızlı 2-opt ve Or-opt uygulamalarının aday listeleri ve bakma
bitleri (don't-look bits) \citep{bentley_1992_fast} ile pertürbasyon sayısının çözüm kalitesindeki
belirleyici rolü \citep{johnson_1997_case} bu çatının pratik temelidir. \citet{applegate_2003_clk}
büyük örneklerde vuruşu turun küçük bir bölgesine sınırlamanın (segment, rastgele yürüyüş,
geometrik vuruşlar) önemini gösterir. LKH \citep{helsgaun_2000_lkh} ve LKH-3 \citep{helsgaun_2017_lkh3}
bu çizginin en güçlü çözücüleridir. Bu çalışmalar vuruşun \emph{nerede} yapılacağını çoğunlukla
rastgele ya da yerellik kuralıyla belirler. Double-bridge seçiminin rastgele yerine yanlı
yapılmasının yararı ILS çatısında daha önce tartışılmıştır \citep{lourenco_2010_ils}; bizim
katkımız, aynı motoru sabit tutup yanlılığın kaynağını öğrenilmiş ve eğitimsiz skorlar arasında eş
bütçede karşılaştırmaktır.

\subsection{Uzun kenar bilgisiyle yönlendirme}
GLS, yerel optimumdaki kenarları $c_e/(1+p_e)$ fayda ölçütüyle cezalandırır: uzun ve az
cezalandırılmış kenarlar önce hedeflenir \citep{voudouris_1999_gls}. Bilgi güdümlü yerel arama (knowledge-guided local search, KGLS) aynı fikri araç rotalamada pertürbasyon aşamasına taşır: bir veri madenciliği çalışmasından türetilen rota genişliği ve kenar uzunluğu ölçütlerine göre ``kötü'' kenarları cezalandırır ve yerel aramayı bu kenarlardan başlatır. Yazarlar bu yönlendirmenin rastgele cezalandırmadan daha iyi sonuç verdiğini, ancak etkisinin yerel arama operatörlerinin seçiminden küçük kaldığını bildirir \citep{arnold_2019_kgls}. Basit
çözüm ve örnek ölçütlerinin iyi ve kötü çözümleri ayırt edebildiği de gösterilmiştir
\citep{arnold_2019_vrpgood}. LKH'nin alfa-yakınlık aday kümeleri ise bir kenarın değerini, o kenar zorunlu kılındığında en küçük 1-ağacın uzunluğundaki artışla ölçer \citep{helsgaun_2000_lkh}. Fazla-uzunluk vekilimiz, GLS'nin uzun kenar sezgisini kenar başına $O(1)$ işlem maliyetiyle
double-bridge pertürbasyon hedeflemesine uyarlar. Katkımız bu yalın uyarlamanın öğrenilmiş
skorlara karşı eş duvar-saatinde ve ön-kayıtlı bir protokolle sınanmasıdır.

\subsection{Öğrenilmiş arama rehberliği}
Uçtan uca inşa yöntemleri turu doğrudan üretir \citep{bello_2017_nco, dai_2017_s2v, kool_2019_attention, kwon_2020_pomo};
iyileştirme yöntemleri hamle seçimini öğrenir \citep{wu_2022_improvement}. Daha güçlü sonuçlar
veren kol, öğrenilmiş bileşeni klasik bir aramanın içine koyar. \citet{hudson_2022_gnngls} GNN'nin
kenar pişmanlık tahminlerini GLS cezalarına bağlar; \citet{sui_2024_neuralgls} bu fikri geliştirir.
NeuroLKH aday kümelerini \citep{xin_2021_neurolkh}, \citet{zheng_2023_rlkh} aynı kararları takviye
öğrenmesiyle öğrenir. Isı haritası yöntemleri kenar olasılıklarını aramaya taşır \citep{fu_2021_generalize, qiu_2022_dimes, sun_2023_difusco}; hiyerarşik ayrıştırma ise öğrenilmiş politikaları çok büyük örneklere ölçekler \citep{pan_2023_htsp}. Taramalar bu yöntemleri öğrenmenin çözücüdeki rolüne (uçtan uca, iyileştirme, klasik sezgiselle melez) ya da metasezgiselin hangi bileşenine (başlatma, operatör seçimi, parametre ayarı) katıldığına göre sınıflar \citep{bengio_2021_mlco, karimimamaghan_2022_mlmeta, sui_2025_survey}; eğitim dağılımı dışına
genellemenin zor olduğu da bilinir \citep{joshi_2020_generalization}. Bizim GNN, üyelik-MLP ve
kazanç-MLP skorlayıcılarımız bu ailenin vuruş hedeflemedeki temsilcileridir. Farkımız, bunları
aynı bilgiden türetilen eğitimsiz vekillerle eş duvar-saatinde karşılaştırmamızdır.

\subsection{Basit tabanlara karşı sınama}
Birkaç çalışma, öğrenilmiş bileşenin katkısını eğitimsiz bir tabana karşı ayrıştırmanın neden
gerekli olduğunu göstermiştir. Maksimum bağımsız küme (maximum independent set, MIS) için GNN güdümlü ağaç aramasında
GNN çıktısı rastgele değerlerle değiştirildiğinde başarım benzer kalmıştır \citep{bother_2022_wrong};
modern GNN'lerin aynı problemde klasik açgözlü algoritmanın gerisinde kaldığı da bildirilmiştir
\citep{angelini_2023_gnngreedy}. Bu sonuçlar MIS'e aittir ve burada yalnız benzetme olarak anılır.
TSP'de \citet{xia_2024_rethinking}, uzaklıktan türetilen basit bir ısı haritasının öğrenilmiş ısı
haritalarıyla rekabet ettiğini gösterir. Rotalama için önerilen hesaplamalı test kılavuzu güçlü
klasik tabanlar ve adil zaman raporu ister \citep{accorsi_2022_guidelines}. Çalışmamız bu deseni
ILS vuruş hedeflemesine, dokunulmamış bir sette ve ön-kayıtlı testlerle taşır.

\subsection{Uyarlanabilir operatör seçiminde kredi ataması}
Reaktif arama, parametreleri arama geçmişinden uyarlar \citep{battiti_2009_reactive}. Uyarlanabilir
operatör seçimi (adaptive operator selection, AOS) çok kollu haydut çerçevesini kullanır
\citep{auer_2002_ucb, dacosta_2008_aos}; uyarlanabilir geniş komşuluk araması da operatörleri geçmiş
başarımlarına göre seçer \citep{ropke_2006_alns}. Bu alanda başarım hem seçim kuralına \citep{thierens_2005_pursuit} hem de operatöre verilen kredinin nasıl tanımlandığına bağlıdır \citep{fialho_2010_aos}. Ortalama ödül nadir ama büyük
iyileşmeleri bastırabilir; uç değerlere dayalı kredi bu nedenle önerilmiştir \citep{fialho_2008_extreme}.
Çevrimiçi vuruş skorlayıcımız bu sınıftandır. Bulgumuz, vuruş başına ortalama ödülün bu motorda
son kaliteyle hizalanmadığını ölçen bir örnektir; çevrimiçi öğrenmenin genel olarak zararlı olduğu
iddiası değildir.

\subsection{Deneysel yöntembilim}
Sezgisel testinin yarışma yerine kontrollü ve hipotez güdümlü bir deney olması gerektiği uzun
süredir savunulur \citep{hooker_1995_testing}. Deneysel algoritma çalışmalarında örnek seçimi,
tekrarlanabilirlik ve raporlama kuralları \citet{johnson_2002_guide}, \citet{mcgeoch_2012_guide} ve
\citet{bartzbeielstein_2020_benchmarking} tarafından derlenmiştir. Ön-kayıt, doğrulayıcı tahmini
sonradan yapılan açıklamadan ayırır \citep{nosek_2018_prereg}. Çok problemli karşılaştırmada örneği
birim alan Wilcoxon testi ve Holm düzeltmesi yerleşik uygulamadır
\citep{demsar_2006_comparisons, wilcoxon_1945, holm_1979}. ``Fark yok'' iddiası ise önceden belirlenmiş
bir sınırla kurulan tek yönlü testler ister \citep{schuirmann_1987_tost, lakens_2017_equivalence}.
Protokolümüz bu öğeleri, geliştirme ve doğrulama setlerini ayıran iki aşamalı bir zincirde birleştirir.

\section{Yöntem}\label{sec:yontem}

Bütün stratejiler aynı ILS motorunu paylaşır ve yalnız double-bridge vuruşunun kesim konumlarını
nasıl seçtiklerinde ayrılır. Hiçbir model kabul kararı vermez: sinyaller yalnız vuruşun nereye
yapılacağını önerir, kabul her zaman gerçek tur maliyetiyle yapılır. Böylece stratejiler
arasındaki her fark doğrudan hedef seçimine bağlanabilir.

\subsection{Problem gösterimi}\label{sec:yontem-problem}

Bir TSP örneği $I=(V,d)$ ile gösterilsin; $n=|V|$ şehir içerir ve $d$ örneğin bildirdiği uzaklık
işlevidir (EUC\_2D'de yuvarlatılmış Öklid uzaklığı, GEO'da TSPLIB büyük çember uzaklığı). Tur
$T=(v_1,\dots,v_n)$, $v_{n+1}=v_1$ olmak üzere döngüsel bir permütasyondur ve maliyeti
$C(T)=\sum_{i=1}^{n} d(v_i,v_{i+1})$'dir. Kalite, bilinen optimuma ya da referans çözüme (sentetik problemlerde LKH-3'ün 60~sn'lik
koşusu) göre yüzde farkla ölçülür: $\mathrm{gap}(T)=100\,(C(T)-C^\star)/C^\star$. İki strateji arasındaki fark gap
\emph{puanı} olarak verilir; negatif değer soldaki stratejinin lehinedir. $N_k(u)$, $u$'nun $k$
en yakın komşusunun listesidir.

\subsection{Hızlı ILS test ortamı}\label{sec:yontem-hizli}

Bir vuruş politikasının etkisi, ancak vuruş döngüsü bütçe içinde binlerce kez dönebiliyorsa
ölçülebilir. Bu amaçla standart hızlı yerel arama tekniklerini \citep{bentley_1992_fast,
johnson_1997_case} izleyen ve Numba ile makine koduna derlenen bir ILS kurduk (\path{fast_ils.py}).
Aynı dosya Aşama~2 ve Aşama~3'te değiştirilmeden kullanıldı; SHA-256 özeti iki protokolde de
kayıtlıdır.

\begin{itemize}
  \item \textbf{Yerel arama.} 2-opt ve Or-opt (1--3 şehirlik segmentin iki yönde yeniden
  yerleştirilmesi); her şehir için $N_{10}$ aday listesi ve $d(a,c)<d(a,\mathrm{succ}(a))$ budama
  kuralı. 2-opt ters çevirmesi kısa olan yay üzerinde yapılır.
  \item \textbf{Bakma bitleri.} Etkin şehirler bir kuyrukta tutulur. Vuruştan sonra yalnız değişen
  kenarların uçları kuyruğa girer; bir vuruşun bedeli tüm tur taraması yerine yerel onarımla
  sınırlı kalır \citep{bentley_1992_fast}.
  \item \textbf{Kabul.} Aday yalnız $C(T')<C(T)$ ise kabul edilir; aksi hâlde tur geri yüklenir.
  Maliyet artımlı izlenir ve koşu sonunda tam yeniden hesapla doğrulanır.
  \item \textbf{Başlangıç.} En yakın komşu turu (başlangıç şehri tohuma bağlı) ve tam 2-opt/Or-opt.
  Bir (problem, tohum) çiftinde bütün stratejiler aynı başlangıç turunu paylaşır; başlangıç süresi
  bütçeye dahil değildir.
\end{itemize}

Amaç LKH-3 ile yarışmak değil, vuruş politikasının gerçekten devrede olduğu, açık ve yeniden
üretilebilir bir ölçüm düzeneği sağlamaktır.

\subsection{Vuruş geometrisi ve havuz kuralı}\label{sec:yontem-kick}

Double-bridge \citep{martin_1992_largestep, lourenco_2010_ils}, kesim konumları $a<b<c$ için turu
$T_{[:a]}\,T_{[b:c]}\,T_{[a:b]}\,T_{[c:]}$ biçiminde yeniden bağlar. \emph{Global} geometride
kesimler $[1,n-1]$ aralığında, birbirinden en az iki konum uzakta seçilir. Aşama~2 ayrıca
\emph{yerel} (segment) geometriyi de sınadı: ilk kesimden sonra diğer ikisi en fazla $W=50$ konum
içinde seçilir \citep{applegate_2003_clk}. Aşama~3 yalnız global geometriyi kullanır.

Hedefli stratejiler her tazelemede bir \emph{havuz} kurar: sinyale göre en şüpheli
\[
  m=\max\bigl(12,\lfloor f\,n\rfloor\bigr)
\]
tur kenarı. Global vuruşun üç kesiminin her biri, $p$ olasılıkla havuzdan tekdüze seçilen bir
kenarın üzerine konur, yani o kenar kırılır; aksi hâlde kesim tekdüzedir. Havuz oranı $f$ ve hedef
olasılığı $p$ havuz kuralını tanımlar. Varsayılan kural $f=0{,}06$, $p=0{,}75$'tir; Aşama~3'te üç
sinyal için kural geliştirme setlerinde seçilmiştir (Bölüm~\ref{sec:deney-asama3}). Tekdüze
kesimler, hedefli stratejilerde de keşfin sürmesini sağlar.

\subsection{Hedefleme sinyalleri}\label{sec:yontem-sinyal}

Sekiz sinyal üç gruba ayrılır: eğitimsiz geometrik vekiller, çevrimdışı eğitilmiş skorlar ve koşu
içinde öğrenen skor. Hepsi aynı havuz kuralını ve aynı maliyet muhasebesini kullanır.

\paragraph{Rastgele.} Havuz yoktur; bütün kesimler tekdüzedir. Bu, ILS'nin klasik davranışıdır ve
ana kontrol koludur.

\paragraph{kNN sırası.} Tur kenarı $(u,v)$ için şüphe skoru
$s_{\mathrm{kNN}}(u,v)=\tfrac12\bigl(r_u(v)+r_v(u)\bigr)$'dir; $r_u(v)$, $v$'nin $N_{24}(u)$
içindeki sırasıdır (listede yoksa $30$). Eşitlikler, kenar uzunluğunun uç noktaların en yakın komşu
uzaklıklarına oranıyla bozulur. Skor hiçbir eğitim ya da çözücü etiketi kullanmaz.

\paragraph{Fazla-uzunluk vekili.} İyi bir turda bir şehre bağlanan kenarlar tipik olarak en yakın
komşularına giden kenarlar ölçeğindedir. Bu ölçekten ne kadar uzun olduğu, bir kenarı kırmanın
olası kazancı için kaba ama ucuz bir göstergedir. Vekil
\begin{equation}
  e(u,v)=d(u,v)-\tfrac12\bigl(d_{\mathrm{ref}}(u)+d_{\mathrm{ref}}(v)\bigr),\qquad
  d_{\mathrm{ref}}(u)=\tfrac12\bigl(d(u,u_{(1)})+d(u,u_{(2)})\bigr)
  \label{eq:excess}
\end{equation}
ile tanımlanır; $u_{(1)}$ ve $u_{(2)}$, $u$'nun birinci ve ikinci en yakın komşusudur. Havuz,
$e$ değeri en büyük $m$ kenardır. $d_{\mathrm{ref}}$ örnek başına bir kez hesaplanır; böylece
tazeleme kenar başına $O(1)$'dir. Vekil, aşağıdaki kazanç hedefinin eğitimsiz karşılığı olarak
tasarlandı ve GLS'nin uzun kenar ilkesini \citep{voudouris_1999_gls} vuruş hedeflemeye taşır. Sıra
tabanlı kNN skorundan farkı, kenarın komşu listesindeki yerine değil, uzunluğun büyüklüğüne
bakmasıdır. Protokol sonrasında, keşifsel olarak, $d_{\mathrm{ref}}$ terimi atılmış ham uzunluk
sinyali ($e(u,v)=d(u,v)$; GLS'nin cezasız fayda ölçütü) de sınandı. Bu kol, fazla-uzunluk için
geliştirmede seçilen havuz kuralını (0,06/0,5) olduğu gibi kullandı. % E751

\paragraph{GNN.} Örnekle ve topla (GraphSAGE) tarzında, saf NumPy ile gerçeklenmiş küçük bir
ileti geçişli ağdır: düğüm başına 10 öznitelik (komşu tur kenarlarının uzunlukları, kaldırma
kazancı, en yakın komşu uzaklıkları, dönüş açısı, segmentten sapma ve tur komşularının kNN
sırası), tur kenarı başına 2 öznitelik (z-normalize uzunluk ve komşu sırası), 3 ileti geçişi
katmanı, 16 boyutlu gizli durum ve $N_{24}$ üzerinde geometrik ileti geçişi. Kenar çıkışı, bir tur
kenarının LKH-3 \citep{helsgaun_2017_lkh3} turunda bulunma olasılığını $\hat p_e$ olarak tahmin eder.
Model $n\in[101,400]$ aralığındaki 120 sentetik problemde (beş dağılım) 954 tur anlık görüntüsüyle
eğitilmiş ve dondurulmuştur. Havuz, $\hat p_e$ değeri en düşük $m$ kenardır. Modelin girdileri
arasında kNN sırası doğrudan bulunduğundan, öğrenmenin bu sinyalin ötesine geçip geçmediği ampirik
bir sorudur.

Bu GNN yayımlanmış bir modelin yeniden üretimi değil, yazarın kendi gerçeklemesidir; LKH turlarıyla
eğitilmiş kenar skorları sınıfının küçük bir temsilcisi olarak kullanılır. İşlevselliği iki ölçümle
gösterilir: yanlış kenarları kNN'den daha iyi ayırt eder (Bölüm~\ref{sec:bulgular-asama2}) ve iki
aşamada da rastgele vuruştan iyidir. % E889, E803, E611
Test problemleri ($n=1.376$--$9.882$) eğitim aralığının çok üstündedir; bu nedenle karşılaştırma
bir boyut genellemesi sınamasını da içerir ve vardığımız sonuçlar bu ölçekteki öğrenilmiş skorlarla
sınırlıdır. % E797, E598

\paragraph{Üyelik-MLP ve kazanç-MLP.} GNN'nin kaybının öğrenmeden mi maliyetten mi geldiğini
ayırmak için, GNN'nin eğitim verisinde iki küçük MLP (17 girdi, 32 ve 16 gizli birim) eğitildi.
Girdiler $N_{24}$ listelerinden kenar başına $O(1)$'de hesaplanır: iki uçtaki karşılıklı sıralar ve
bunların en küçük/en büyük değeri, kenar uzunluğunun 1., 5. ve 24. komşu uzaklığına oranları,
komşu iki tur kenarının uzunluk oranları ve sıraları. \emph{Üyelik-MLP} kenarın referans turda
bulunmamasını sınıflandırır. \emph{Kazanç-MLP} ise
\[
  y_e=\mathbb{1}[e\notin T^\star]\cdot\max\bigl(0,\,d_e-\bar d^{\,\star}_e\bigr)/\tilde d
\]
hedefini $\log(1+y_e)$ üzerinden regresyonla öğrenir; $\bar d^{\,\star}_e$, $e$'nin iki ucundaki
referans tur kenarlarının ortalama uzunluğu, $\tilde d$ medyan tur kenarı uzunluğudur. Kazanç
hedefi, vuruşun yanlış kenarı kırmasının yanında bunun ne kadar kazandırabileceğini de ölçer;
Denklem~\eqref{eq:excess}, $\bar d^{\,\star}$ yerine eğitimsiz $d_{\mathrm{ref}}$ koyarak bu hedefe
yaklaşır.

\paragraph{Sabit doğrusal ön-bilgi.} Aynı 17 öznitelik ve bir sabit terim üzerinde, kazanç hedefine
ridge regresyonuyla uydurulmuş doğrusal bir skor $w_0^\top x_e$'dir. Koşu boyunca değişmez; havuz en
yüksek skorlu $m$ kenardır. Bu kol, çevrimiçi öğrenmenin etkisini modelden ayırmak için eşleşik
kontroldür.

\paragraph{Çevrimiçi (Thompson-ridge).} Aynı doğrusal model koşu sırasında, her vuruşun gerçekleşen
sonucundan öğrenir. Kesimleri $e_1,e_2,e_3$ olan bir vuruşun ödülü, yerel arama sonrasındaki
maliyet değişiminden
\begin{equation}
  r=-\frac{C(T_{\text{sonra}})-C(T_{\text{önce}})}{\tilde d_0}
  \label{eq:reward}
\end{equation}
olarak hesaplanır; $\tilde d_0$ başlangıç turunun medyan kenar uzunluğudur. İyileşme pozitif ödül
verir. Ödül toplamsal modellenir, $r\approx w^\top(x_{e_1}+x_{e_2}+x_{e_3})$, ve $w$ çevrimiçi
ridge regresyonuyla güncellenir: $A=\lambda I+\sum xx^\top$, $b=\lambda w_0+\sum r\,x$,
$\lambda=50$. Her tazelemede $w\sim\mathcal N(A^{-1}b,\ \nu^2A^{-1})$, $\nu=0{,}5$, örneklenir ve
tahmini ödülü en yüksek $m$ kenar havuzu oluşturur. Model $w_0$ ile ılık başlatılır. Havuz, diğer sinyallerdeki genel kuralla
(Bölüm~\ref{sec:yontem-maliyet}) her kabul edilen vuruşta, ayrıca 2.000 vuruşta bir yeniden kurulur. İşaretli ödülün bilinçli bir gerekçesi
vardır: vuruşların çok küçük bir kısmı kabul edildiğinden, $\max(0,\cdot)$ biçimindeki bir ödül
neredeyse her vuruş için sıfır olurdu. Ödül işareti ve havuz seçimi birim testleriyle doğrulandı:
iyileştiren vuruş pozitif, nötr vuruş sıfır ödül alır ve havuz en yüksek tahmini ödülü seçer. % E959
Aşama~3 geliştirme koşularında soğuk başlatılan ($w_0=0$) bir sürüm de sınandı.

\subsection{Maliyet muhasebesi}\label{sec:yontem-maliyet}

Havuz, incumbent tur değiştiğinde yeniden hesaplanır. Skorlama süresi (GNN çıkarımı, MLP ileri
geçişi, vekil hesabı, Thompson örneklemesi) duvar-saati bütçesinden düşülür; yani her sinyal
ürettiği bilginin bedelini aynı bütçeden öder. Toplam skorlama süresi geçen sürenin \%10'unu
aşarsa tazeleme ertelenir. Bu sırada havuzdaki bir kenar artık turda değilse o kesim tekdüze
seçilir; buna \emph{bayat kesim} diyoruz. Örnek başına bir kez yapılan hazırlıklar ($N_{24}$
listeleri, komşu uzaklıkları, ödül ölçeği) hiçbir sinyal için bütçeye yüklenmez. Numba derleme
süresi her işçi sürecinde bir ısınma çağrısıyla bütçe dışına alınır.

\subsection{Aşama 1 iskeleti}\label{sec:yontem-crine}

Aşama~1, öğrenilmiş vuruş hedeflemesini CRINE-ILS (Counterfactual Residual Improvement through Neighborhood Expansion) adlı yavaş bir iskelette sınamıştı: saf Python
yerel arama, her vuruştan sonra tüm tur üzerinde sınırlı onarım ve ikinci bir GNN ile sıralanan
şehirler çevresinde komşuluk genişletme (\path{adaptive_ils.py}). Bu iskelette bir vuruş yüzlerce
milisaniye sürüyordu; kısa bütçede vuruş döngüsü neredeyse hiç çalışmadığı için stratejiler
arasındaki farklar tohum gürültüsüne indirgendi (Bölüm~\ref{sec:bulgular}). Hızlı test ortamı bu
dersin doğrudan sonucudur.

\subsection{İşlem karmaşıklığı}\label{sec:yontem-karmasiklik}

Bir tazelemede fazla-uzunluk vekili, kNN skoru ve iki MLP kenar başına sabit iş yapar; toplam
$O(n)$'dir. GNN çıkarımı da $k=24$ ve gizli boyut sabit olduğundan $O(n)$'dir, ancak sabit çarpanı
birkaç ileti geçişi katmanı nedeniyle büyüktür. Çevrimiçi skorun her vuruştaki güncellemesi 18
boyutlu bir dış çarpım, her tazelemedeki örneklemesi $18\times18$ bir ters alma gerektirir. Global
vuruş ve reddedilen vuruştan sonra tur geri yüklemesi $O(n)$ dizi kopyasıdır; vuruş sonrası onarım,
etkinleşen şehir sayısıyla orantılıdır.

\section{Deneysel kurulum}\label{sec:deney-kurulum}

\subsection{Aşama 1 ve Aşama 2}\label{sec:deney-onkayit}

Aşama~2'nin bütün kararları ilk koşudan önce \path{PROTOCOL.json} dosyasında sabitlendi ve dosyanın
SHA-256 özeti kaydedildi (\texttt{cce83992\ldots}). % E791
Birincil set, projede daha önce kullanılmamış 38 VLSI problemidir (VLSI-38; $n=1.376$--$3.954$,
10 tohum, 10~sn). % E797
Dört birincil hipotez GNN'yi rastgele ve kNN hedeflemeyle iki geometride karşılaştırır: Wilcoxon
testleri tek Holm ailesinde düzeltilir, ``fark yok'' sonucu $\pm0{,}10$ puanlık sınırla iki tek
yönlü testle (TOST) sınanır. % E792
Tablo~\ref{tab:v2-sets} ikincil setleri verir. Bellek baskısıyla iki kez kesilen koşu, tamamlanmış
kayıtlara dokunulmadan sürdürüldü.

% OTOMATIK URETIM -- evidence/scripts/build_v2.py
\begin{table}[!htbp]
\centering\small\setlength{\tabcolsep}{3pt}
\caption{Hızlı ILS kampanyasının deney setleri. Her (problem, tohum) için altı strateji aynı başlangıç turundan ve aynı kick rastgele sayı tohumundan başlar.}
\label{tab:v2-sets}
\begin{tabularx}{\textwidth}{@{}>{\raggedright\arraybackslash}Xcccccl@{}}
\toprule
Set & Problem & $n$ & Tohum & Bütçe [sn] & Referans & Rol \\
\midrule
VLSI-38 (birincil) & 38 & 1376--3954 & 10 & 10 & Optimum/BKS & Ön-kayıtlı birincil \\
VLSI-20 (önceki) & 20 & 131--3293 & 10 & 10 & Optimum/BKS & İkincil \\
TSPLIB-72 & 72 & 51--2392 & 5 & 10 & Optimum/BKS & İkincil \\
Sentetik-100 (1 sn) & 100 & 81--160 & 5 & 1 & LKH-3 (30 sn) & İkincil \\
Sentetik-1000 & 40 & 1000--1000 & 5 & 10 & LKH-3 (30 sn) & İkincil \\
\bottomrule
\end{tabularx}
\end{table}


Birincil sonuçtan sonra eklenen kâhin hedefleme, sürekli tazelenen GNN, 60~sn bütçe ve iki MLP
deneyleri keşifseldir.

\subsection{Aşama 3: geliştirme ve seçim}\label{sec:deney-asama3}

Aşama~3'ün yöntem ve ayar seçimleri yalnız geliştirme setlerinde yapıldı: VLSI-38, Aşama~1'in 20
VLSI problemi (VLSI-20) ve Sentetik-1000 (5 tohum, 10~sn). Bu koşularda kNN ve fazla-uzunluk
VLSI-38'de rastgeleden iyiydi ve birbirinden ayrışmadı. % E738, E739, E740
Çevrimiçi model VLSI-38'de eşleşik sabit kolundan ($+0{,}210$ puan), diğer iki sette fazla-uzunluk
vekilinden kötüydü ($+0{,}157$ ve $+0{,}148$; ikincisi soğuk başlatılmış sürüm). % E742, E745, E747
Ardından kNN, fazla-uzunluk ve GNN için $f\in\{0{,}02;0{,}04;0{,}06;0{,}10\}$ ve
$p\in\{0{,}5;0{,}75;0{,}9\}$ ızgarası VLSI-38 ile Sentetik-1000'de 3 tohumla koşuldu. Seçim kuralı
ızgara sonuçlarından önce sapma kaydına yazıldı: iki setin ortalama gap'ini en küçükleyen hücre
seçilir; varsayılan kural en iyiye 0,01 puandan yakınsa korunur. Kural kNN için 0,04/0,75,
fazla-uzunluk için 0,06/0,5, GNN için 0,04/0,5'i seçti (Şekil~\ref{fig:v3-pool-grid}). % E750, E751, E752
Kazanç-MLP dahil diğer sinyaller ayarlanmadı.

\subsection{Aşama 3: dondurma, zaman damgası ve doğrulama seti}\label{sec:deney-dondurma}

Protokol \path{PROTOCOL_v2.json} dosyasında donduruldu; dosya hipotezleri, analizi ve kod, model ve
veri dosyalarının SHA-256 özetlerini içerir, kendi özeti \texttt{96eb1d9e\ldots8085}'tir. % E584
Kod, doğrulama setini protokol donmadan çalıştırmayı reddeder. Protokolün öngördüğü Open Science
Framework yazar için erişilemez olduğundan bu sapma kayda yazıldı; özet 12:51'de
OpenTimestamps\footnote{\url{https://opentimestamps.org}} takvim sunucularına gönderildi ve holdout2
koşusu hemen ardından başladı. Bitcoin tasdiki (blok 976258), koşu 540 görevin 384'ündeyken ve
sonuçlara bakılmadan alındı. % E585, E964
Bu blok zamanı, protokolün sonuçlar analiz edilmeden önce var olduğunu ve sonradan
değiştirilmediğini gösteren, üçüncü kişilerce bağımsız doğrulanabilen ve kurcalamaya dayanıklı üst
sınırdır; koşudan önceki sıralama yerel sapma kaydına dayanır.

Holdout2, 54 problemden oluşur: 6 Waterloo VLSI, 3 National ve 3 TSPLIB problemi ile 14 uniform,
14 kümeli ve 14 karışık sentetik problem ($n=1.689$--$9.882$). % E586, E591--E598
TSPLIB problemleri Aşama~1'de kullanılmış, ancak hiçbir Aşama~2 ya da Aşama~3 koşusunda yer
almamıştır. Sekiz strateji 10 tohum ve 10~sn ile koşuldu (toplam 4.320 koşu). % E587, E588, E590
Bir (problem, tohum) çiftinin stratejileri aynı işçide ardışık çalışır ve vuruş tohumunu da
paylaşır. Dört işçi Intel Core i5-9300H, Windows~11 x64, Python~3.13 ve Numba~0.68
ile çalıştı; bütçe duvar-saati olduğundan vuruş sayıları donanıma bağlıdır
(Tablo~\ref{tab:v3-mechanism}). Hiçbir koşu dışlanmadı.

\subsection{Aşama 3: hipotezler ve istatistik}\label{sec:deney-istatistik}

İstatistik birimi problemdir: gap önce 10 tohum üzerinden ortalanır ($n=54$). Dört doğrulayıcı
hipotez sabit sırayla, her biri ancak öncekiler reddedildiyse $\alpha=0{,}05$ düzeyinde sınanır.
V1 çevrimiçi skorun sabit doğrusal ön-bilgiden kötü, V2 vekilin rastgeleden iyi olduğunu öne sürer
(iki taraflı Wilcoxon \citep{wilcoxon_1945}). V3 ve V4, vekilin GNN'ye ve kazanç-MLP'ye
$\delta=0{,}05$ puandan fazla kötü olmadığını $d-\delta$ üzerinde tek yönlü Wilcoxon testiyle sınar
\citep{schuirmann_1987_tost, lakens_2017_equivalence}. % E589
Sınır, geliştirmede gözlenen en küçük anlamlı hedefleme etkisinin (fazla-uzunluk$-$rastgele,
$-0{,}070$ puan) altındadır; daha dar bir sınır, problem başı farkların standart sapması 0,131'e
ulaştığı için güçsüz kalırdı. % E739, E740
Öte yandan sınır, GNN'nin rastgeleye karşı Aşama~2'deki etkisiyle ($-0{,}052$ puan) aynı büyüklüktedir;
bu yüzden V3 tek başına GNN kazancının korunduğunu göstermez, V2 ve keşifsel üstünlük
karşılaştırmasıyla birlikte okunmalıdır. % E803
İkincil aile beş keşifsel üstünlük karşılaştırmasıdır (kNN$-$rastgele, fazla-uzunluk$-$kNN,
GNN$-$kNN, kazanç-MLP$-$üyelik-MLP, fazla-uzunluk$-$GNN) ve Holm ile düzeltilir \citep{holm_1979}.
Kayıtlı analiz betiği `:' içeren strateji adlarını ayrıştıramadığından aile, aynı istatistik
işlevlerini çağıran ayrı bir betikle (\path{build_v3.py}) hesaplandı. Güven aralıkları 10.000
tekrarlı bootstrap'la bulunur.

\section{Bulgular}\label{sec:bulgular}

Ana sonuçlar, Aşama~2 ve 3'ün hiçbir geliştirme koşusunda kullanılmamış 54 problemden (holdout2),
10 tohum ve 10~sn bütçeyle elde edildi. Farklar gap puanıdır; negatif değer ilk stratejinin lehinedir.
Bölüm önce ön-kayıtlı doğrulayıcı testi, ardından keşifsel katman, mekanizma, duyarlılık ve
sağlamlık analizlerini, sonra yayımlanmış bir öğrenilmiş modelle, eşit vuruş sayısıyla ve LKH-3 ile
yapılan hakem odaklı blokları verir. Son alt bölüm, Aşama~3 hipotezlerini doğuran Aşama~2 sonuçlarını
özetler: Aşama~2'de eğitimsiz kNN sırasının GNN'yi geçmesi, eğitimsiz bir vekilin yeterli olabileceği
fikrini doğurdu; Aşama~3 bu fikri yeni bir test setinde ön-kayıtlı olarak sınadı.
% kaynak: E586-E588

\subsection{Doğrulayıcı test}\label{sec:bulgular-dogrulayici}

Sabit sırayla sınanan dört ön-kayıtlı hipotezin dördü de desteklendi
(Tablo~\ref{tab:v3-confirm}). Koşu içinde ağırlıklarını güncelleyen çevrimiçi model, aynı ağırlıklarla
başlayıp hiç güncellenmeyen sabit doğrusal modelden $0{,}291$ puan kötüdür ve 54 problemin hiçbirinde
onu geçemez (V1, $p=1{,}6\times10^{-10}$).
Eğitimsiz fazla-uzunluk vekili, rastgele vuruşa göre ortalama gap'i $0{,}172$ puan düşürür
(V2, $p=3{,}7\times10^{-7}$).
Vekil, LKH-3 turlarıyla eğitilmiş GNN'ye (V3) ve kazanç hedefiyle eğitilmiş MLP'ye (V4) karşı
$0{,}05$ puanlık sınırla kötü-değil testlerini geçer; gözlenen farklar vekil lehinedir ($-0{,}119$
ve $-0{,}083$ puan; $p=7{,}9\times10^{-9}$ ve $1{,}8\times10^{-7}$).
Problem başı farklar sağa çarpık olduğundan Hodges--Lehmann kestirimleri ortalamalardan küçüktür
(V1 $+0{,}242$ [$+0{,}167$; $+0{,}337$], V2 $-0{,}101$ [$-0{,}166$; $-0{,}062$], V3 $-0{,}069$
[$-0{,}114$; $-0{,}039$], V4 $-0{,}057$ [$-0{,}082$; $-0{,}032$]); kararlar değişmez, V3
ve V4'te üst sınırlar ($-0{,}039$ ve $-0{,}032$) sıfırın altındadır.
Aşama~1'de kullanılmış üç TSPLIB problemi çıkarıldığında da (51 problem; keşifsel duyarlılık
denetimi) dört hipotezin dördü reddedilir (V1 $+0{,}295$, V2 $-0{,}166$, V3 $-0{,}114$, V4 $-0{,}084$
puan; en büyük $p=7{,}5\times10^{-7}$).
Ön-kayıtlı iddia bu dört sonuçtan ibarettir. Aynı tablodaki Holm düzeltmeli ikincil çiftler
keşifseldir: vekil GNN'den ($p_{\mathrm{Holm}}=1{,}0\times10^{-4}$) ve kNN sırasından ($-0{,}093$ puan,
$p_{\mathrm{Holm}}=3{,}4\times10^{-5}$) iyidir; GNN ile kNN ayrışmaz ($p_{\mathrm{Holm}}=1{,}0$).
% kaynak: E589, E599-E602, E604, E605, E607, E960-E963, E1100-E1103

\begin{table}[!htbp]
\centering\small\setlength{\tabcolsep}{3pt}
\caption{Doğrulayıcı test (PROTOCOL\_v2, holdout2: 54 problem, 10 tohum, 10~sn). $\Delta$: A$-$B gap farkı [puan], negatif A lehine; kötü-değil sınırı 0{,}05. Renk: anlamlı fark. Ö/E/G: A'nın önde, eşit ve geride olduğu problem sayısı.}
\label{tab:v3-confirm}
\begin{tabularx}{\textwidth}{@{}l>{\raggedright\arraybackslash}Xlccccl@{}}
\toprule
ID & Karşılaştırma (A$-$B) & Test & $\Delta$ & \%95 GA & Ö/E/G & $p$ & Karar \\
\midrule
V1 & Çevrimiçi $-$ Sabit doğrusal & üstünlük & \kotu{+0{,}291} & [+0{,}221; +0{,}367] & 0/0/54 & $1{,}6\times10^{-10}$ & $H_0$ ret \\
V2 & Fazla-uzunluk $-$ Rastgele & üstünlük & \iyi{$-$0{,}172} & [$-$0{,}255; $-$0{,}098] & 44/0/10 & $3{,}7\times10^{-7}$ & $H_0$ ret \\
V3 & Fazla-uzunluk $-$ GNN & kötü-değil & \iyi{$-$0{,}119} & [$-$0{,}195; $-$0{,}051] & 41/0/13 & $7{,}9\times10^{-9}$ & $H_0$ ret \\
V4 & Fazla-uzunluk $-$ Kazanç-MLP & kötü-değil & \iyi{$-$0{,}083} & [$-$0{,}146; $-$0{,}026] & 42/0/12 & $1{,}8\times10^{-7}$ & $H_0$ ret \\
\midrule
\multicolumn{8}{l}{\textit{İkincil, keşifsel (protokolde adı geçen 5 çift; $p$ Holm düzeltmeli)}} \\
-- & kNN $-$ Rastgele & üstünlük & $-$0{,}079 & [$-$0{,}149; $-$0{,}018] & 31/0/23 & 0{,}169 & anlamsız \\
-- & Fazla-uzunluk $-$ kNN & üstünlük & \iyi{$-$0{,}093} & [$-$0{,}144; $-$0{,}040] & 43/0/11 & $3{,}4\times10^{-5}$ & anlamlı \\
-- & GNN $-$ kNN & üstünlük & +0{,}026 & [$-$0{,}027; +0{,}089] & 30/0/24 & 1{,}000 & anlamsız \\
-- & Kazanç-MLP $-$ Üyelik-MLP & üstünlük & $-$0{,}006 & [$-$0{,}052; +0{,}038] & 25/0/29 & 1{,}000 & anlamsız \\
-- & Fazla-uzunluk $-$ GNN & üstünlük & \iyi{$-$0{,}119} & [$-$0{,}195; $-$0{,}051] & 41/0/13 & $1{,}0\times10^{-4}$ & anlamlı \\
\bottomrule
\end{tabularx}
\end{table}


Şekil~\ref{fig:manset}'teki problem düzeyi dağılımı çarpıktır: vekil$-$GNN farkının medyanı
$-0{,}060$ puandır. Ham kazanma oranı 41/54'tür; ölçüm gürültüsünün (0,016 puan) altındaki farklar
beraberlik sayıldığında oran \%70'e iner. Şekilde işaretli örnek problemler, GNN'nin kendi
çeyrekler arası aralığında kaldığı problemler arasından katman başına en büyük avantajla seçildi.
% kaynak: E789, V013, V014

\subsection{Strateji sıralaması ve problem türleri}\label{sec:bulgular-katman}

Tablo~\ref{tab:v3-strategies} sekiz stratejiyi dört problem türünde karşılaştırır. Fazla-uzunluk,
tüm sette en düşük ortalama gap'i verir (\%0,922; rastgele \%1,094, GNN \%1,041) ve dört katmanın
hepsinde sıralamanın başındadır.
Katman sonuçları betimseldir; tablodaki renkler katman içi Holm düzeltmesine dayanır. Hedeflemenin getirisi problem
yapısına bağlıdır: rastgeleye göre kazanç gerçek problemlerde $0{,}323$, kümeli problemlerde
$0{,}302$ puandır; karışık katmanda $0{,}052$ puana iner, uniform katmanda $0{,}030$ puanda kalır ve
anlamlı değildir ($p=0{,}12$).
kNN sırası bu ayrımda kararsızdır. Kümeli problemlerde rastgeleden $0{,}271$ puan iyi, karışık
problemlerde $0{,}039$ puan kötüdür (düzeltmesiz, betimsel) ve tüm sette rastgeleden anlamlı ayrışmaz
($p_{\mathrm{Holm}}=0{,}17$).
Öğrenilmiş üç skor (GNN, kazanç-MLP, üyelik-MLP) tüm sette rastgeleden iyidir ama vekilin gerisinde
kalır; çevrimiçi model her katmanda en kötüdür.
% kaynak: E603, E608, E610-E613, E615, E618, E626, E633, E634, E641, E642

Yapıya bağımlılık, katman etiketinden bağımsız bir ölçüyle de görülür (keşifsel). Her problem için
en yakın komşu uzaklığının değişim katsayısını (CV) bir kümelenme göstergesi olarak kullandık. Bu
değer uniform problemlerde $0{,}51$ ile $0{,}55$ arasındadır, kümeli ve ulusal problemlerde daha
büyüktür.
Vekilin rastgeleye göre kazancı CV arttıkça büyür (Spearman $\rho=-0{,}56$, $p=9{,}8\times10^{-6}$;
Şekil~\ref{fig:v4-structure}). GNN'ye karşı farkta ilişki aynı yöndedir ama daha zayıftır
($\rho=-0{,}38$, $p=0{,}005$); problem boyutuyla ilişki yoktur ($\rho=-0{,}10$, $p=0{,}46$).
Ölçü, doğrulayıcı koşunun mevcut verisi üzerinde hesaplandığı için betimseldir. Uygulayıcı için
sonucu şudur: hedeflemeden en çok, kenar uzunluklarının düzensiz olduğu kümeli ve gerçek problemlerde
yararlanılır; bu ölçü koşudan önce, koordinatlardan ucuzca hesaplanabilir.
% kaynak: E1038-E1041, E1043, E1045

\begin{figure}[!htbp]
\centering
% Problem yapisi ve hedefleme kazanci -- OTOMATIK URETIM: evidence/scripts/build_v4.py
% Veri: figures/data/v4_structure.dat (v4_robustness_analysis.json T_structure.rows)
\begin{tikzpicture}
\begin{axis}[width=\textwidth, height=5.6cm, xmode=log, log basis x=10, xmin=0.25, xmax=2.5, ymax=0.9,
  xtick={0.3,0.5,0.75,1,1.5,2}, xticklabels={0.3,0.5,0.75,1,1.5,2},
  xlabel={Clustering index: CV of nearest-neighbour distance (log scale)},
  ylabel={Excess $-$ random [gap points]}, tick label style={font=\scriptsize}, label style={font=\small},
  ymajorgrids, grid style={black!10},
  legend style={font=\scriptsize, draw=black!40, fill=white, at={(0.01,0.03)}, anchor=south west, legend columns=4}]
\addplot[black!50, thin, forget plot] coordinates {(0.25,0) (2.5,0)};
\addplot[only marks, mark=x, mark size=1.8pt, black!45] table[x=cv, y=d_rand, restrict expr to domain={\thisrow{guni}}{1:1}] {figures/data/v4_structure.dat};
\addlegendentry{Uniform}
\addplot[only marks, mark=square*, mark size=1.3pt, black!70] table[x=cv, y=d_rand, restrict expr to domain={\thisrow{gmix}}{1:1}] {figures/data/v4_structure.dat};
\addlegendentry{Mixed}
\addplot[only marks, mark=*, mark size=1.5pt, iyiRenk] table[x=cv, y=d_rand, restrict expr to domain={\thisrow{gclu}}{1:1}] {figures/data/v4_structure.dat};
\addlegendentry{Clustered}
\addplot[only marks, mark=triangle*, mark size=1.9pt, kotuRenk] table[x=cv, y=d_rand, restrict expr to domain={\thisrow{greal}}{1:1}] {figures/data/v4_structure.dat};
\addlegendentry{Real}
\node[anchor=north west, font=\scriptsize, fill=white, inner sep=2pt] at (rel axis cs:0.01,0.97) {Spearman $\rho=-0.56$, $p=9.8\times10^{-6}$};
\end{axis}
\end{tikzpicture}

\caption{Problem yapısı ve hedefleme kazancı (holdout2, 10~sn, keşifsel): yatay eksen en yakın komşu uzaklığının değişim katsayısı, dikey eksen fazla-uzunluk vekilinin rastgeleye göre problem başı gap farkı (negatif = vekil iyi).}
\label{fig:v4-structure}
\end{figure}

% OTOMATIK URETIM -- evidence/scripts/build_v3.py (elle duzenleme)
\begin{table}[!htbp]
\centering\small\setlength{\tabcolsep}{3pt}
\caption{Holdout2'de sekiz strateji, katman bazında: ortalama gap ve rastgeleye eşlenik fark $\Delta$. Renk: gap sütununda katmanın en iyi/en kötüsü, $\Delta$'da katman içi Holm düzeltmeli $p<0{,}05$; $p$ sütunu düzeltmesizdir.}
\label{tab:v3-strategies}
\begin{tabular*}{\textwidth}{@{\extracolsep{\fill}}lccccc@{}}
\toprule
Strateji & Gap [\%] & $\Delta$ [puan] & \%95 GA & K/B/K & $p$ \\
\midrule
\textit{Tümü (54 problem)}: Rastgele & 1{,}094 & \multicolumn{4}{c}{referans} \\ % E608
\quad kNN (0,04/0,75) & 1{,}015 & $-$0{,}079 & [$-$0{,}149; $-$0{,}018] & 31/0/23 & 0{,}056 \\ % E609
\quad Fazla-uzunluk (0,06/0,5) & \iyi{0{,}922} & \iyi{$-$0{,}172} & [$-$0{,}255; $-$0{,}098] & 44/0/10 & $3{,}7\times10^{-7}$ \\ % E610
\quad GNN (0,04/0,5) & 1{,}041 & \iyi{$-$0{,}053} & [$-$0{,}080; $-$0{,}026] & 37/0/17 & $6{,}8\times10^{-4}$ \\ % E611
\quad Kazanç-MLP & 1{,}005 & \iyi{$-$0{,}089} & [$-$0{,}152; $-$0{,}037] & 36/0/18 & 0{,}003 \\ % E612
\quad Üyelik-MLP & 1{,}011 & \iyi{$-$0{,}083} & [$-$0{,}127; $-$0{,}044] & 35/0/19 & $6{,}6\times10^{-4}$ \\ % E613
\quad Sabit doğrusal & 1{,}012 & \iyi{$-$0{,}082} & [$-$0{,}157; $-$0{,}026] & 34/0/20 & 0{,}020 \\ % E614
\quad Çevrimiçi (ılık) & \kotu{1{,}303} & \kotu{+0{,}209} & [+0{,}150; +0{,}271] & 4/0/50 & $1{,}4\times10^{-8}$ \\ % E615
\addlinespace
\textit{Gerçek (12 problem)}: Rastgele & 1{,}829 & \multicolumn{4}{c}{referans} \\ % E616
\quad kNN (0,04/0,75) & 1{,}724 & $-$0{,}104 & [$-$0{,}280; +0{,}048] & 9/0/3 & 0{,}151 \\ % E617
\quad Fazla-uzunluk (0,06/0,5) & \iyi{1{,}505} & \iyi{$-$0{,}323} & [$-$0{,}506; $-$0{,}159] & 10/0/2 & 0{,}002 \\ % E618
\quad GNN (0,04/0,5) & 1{,}742 & $-$0{,}086 & [$-$0{,}170; $-$0{,}011] & 8/0/4 & 0{,}077 \\ % E619
\quad Kazanç-MLP & 1{,}715 & $-$0{,}114 & [$-$0{,}240; $-$0{,}006] & 9/0/3 & 0{,}092 \\ % E620
\quad Üyelik-MLP & 1{,}704 & $-$0{,}125 & [$-$0{,}213; $-$0{,}034] & 9/0/3 & 0{,}042 \\ % E621
\quad Sabit doğrusal & 1{,}754 & $-$0{,}075 & [$-$0{,}170; +0{,}010] & 8/0/4 & 0{,}233 \\ % E622
\quad Çevrimiçi (ılık) & \kotu{2{,}157} & \kotu{+0{,}328} & [+0{,}208; +0{,}460] & 1/0/11 & $9{,}8\times10^{-4}$ \\ % E623
\addlinespace
\textit{Uniform (14 problem)}: Rastgele & 0{,}606 & \multicolumn{4}{c}{referans} \\ % E624
\quad kNN (0,04/0,75) & 0{,}625 & +0{,}019 & [$-$0{,}005; +0{,}044] & 6/0/8 & 0{,}241 \\ % E625
\quad Fazla-uzunluk (0,06/0,5) & \iyi{0{,}576} & $-$0{,}030 & [$-$0{,}057; $-$0{,}003] & 10/0/4 & 0{,}119 \\ % E626
\quad GNN (0,04/0,5) & 0{,}590 & $-$0{,}016 & [$-$0{,}050; +0{,}019] & 8/0/6 & 0{,}326 \\ % E627
\quad Kazanç-MLP & 0{,}607 & +0{,}001 & [$-$0{,}028; +0{,}027] & 7/0/7 & 0{,}715 \\ % E628
\quad Üyelik-MLP & 0{,}593 & $-$0{,}013 & [$-$0{,}041; +0{,}016] & 8/0/6 & 0{,}502 \\ % E629
\quad Sabit doğrusal & 0{,}592 & $-$0{,}014 & [$-$0{,}045; +0{,}018] & 9/0/5 & 0{,}463 \\ % E630
\quad Çevrimiçi (ılık) & \kotu{0{,}712} & \kotu{+0{,}106} & [+0{,}082; +0{,}131] & 0/0/14 & $1{,}2\times10^{-4}$ \\ % E631
\addlinespace
\textit{Kümeli (14 problem)}: Rastgele & 1{,}400 & \multicolumn{4}{c}{referans} \\ % E632
\quad kNN (0,04/0,75) & 1{,}129 & \iyi{$-$0{,}271} & [$-$0{,}452; $-$0{,}131] & 13/0/1 & $8{,}5\times10^{-4}$ \\ % E633
\quad Fazla-uzunluk (0,06/0,5) & \iyi{1{,}098} & \iyi{$-$0{,}302} & [$-$0{,}533; $-$0{,}097] & 12/0/2 & 0{,}009 \\ % E634
\quad GNN (0,04/0,5) & 1{,}299 & \iyi{$-$0{,}101} & [$-$0{,}146; $-$0{,}053] & 12/0/2 & 0{,}005 \\ % E635
\quad Kazanç-MLP & 1{,}156 & \iyi{$-$0{,}245} & [$-$0{,}429; $-$0{,}110] & 12/0/2 & 0{,}001 \\ % E636
\quad Üyelik-MLP & 1{,}207 & \iyi{$-$0{,}193} & [$-$0{,}306; $-$0{,}092] & 11/0/3 & 0{,}004 \\ % E637
\quad Sabit doğrusal & 1{,}162 & \iyi{$-$0{,}239} & [$-$0{,}482; $-$0{,}065] & 10/0/4 & 0{,}020 \\ % E638
\quad Çevrimiçi (ılık) & \kotu{1{,}725} & \kotu{+0{,}325} & [+0{,}150; +0{,}485] & 3/0/11 & 0{,}005 \\ % E639
\addlinespace
\textit{Karışık (14 problem)}: Rastgele & 0{,}645 & \multicolumn{4}{c}{referans} \\ % E640
\quad kNN (0,04/0,75) & 0{,}683 & +0{,}039 & [+0{,}005; +0{,}070] & 3/0/11 & 0{,}042 \\ % E641
\quad Fazla-uzunluk (0,06/0,5) & \iyi{0{,}592} & \iyi{$-$0{,}052} & [$-$0{,}080; $-$0{,}024] & 12/0/2 & 0{,}004 \\ % E642
\quad GNN (0,04/0,5) & 0{,}632 & $-$0{,}013 & [$-$0{,}042; +0{,}016] & 9/0/5 & 0{,}502 \\ % E643
\quad Kazanç-MLP & 0{,}643 & $-$0{,}002 & [$-$0{,}037; +0{,}031] & 8/0/6 & 0{,}952 \\ % E644
\quad Üyelik-MLP & 0{,}637 & $-$0{,}008 & [$-$0{,}031; +0{,}018] & 7/0/7 & 0{,}463 \\ % E645
\quad Sabit doğrusal & 0{,}645 & +0{,}001 & [$-$0{,}034; +0{,}034] & 7/0/7 & 0{,}903 \\ % E646
\quad Çevrimiçi (ılık) & \kotu{0{,}741} & \kotu{+0{,}096} & [+0{,}070; +0{,}124] & 0/0/14 & $1{,}2\times10^{-4}$ \\ % E647
\bottomrule
\end{tabular*}
\end{table}


\paragraph{Vaka incelemesi.}
Şekil~\ref{fig:manset}'te işaretli kümeli h2c02 probleminde GNN hedeflemesi 10~sn sonunda \%1,47 gap'te kalır, fazla-uzunluk \%0,29'a iner.
GNN burada kendi çeyrekler arası aralığındadır; fark rakibin çöküşünden
gelmez. Gerçek eg7146
probleminde gap $2{,}35$'ten $1{,}53$'e, uniform h2u12'de yalnızca $0{,}58$'den $0{,}46$'ya
düşer; yapısal düzensizlik azaldıkça avantaj küçülür.
Karşı-örnek de kümelidir: h2c13'te GNN \%1,53, vekil \%2,16 gap verir.
Kümeli katmanda vekilin GNN'ye karşı farkının güven aralığı sıfırı keser ($-0{,}201$ puan, GA
$[-0{,}439;\,+0{,}014]$); büyük kazanç, büyük varyansla birlikte gelir.
% kaynak: E671, V001-V008

\subsection{Mekanizma}\label{sec:bulgular-mekanizma}

\paragraph{Ucuz ve taze hedef.}
Tablo~\ref{tab:v3-mechanism}, üstünlüğün verimlilikle uyumlu olduğunu gösterir.
Fazla-uzunluk skorunun duvar-saatindeki payı \%0,65'tir; bu pay GNN'de \%8,77, kNN'de
\%2,83'tür.
GNN havuzu seyrek tazelendiği için kesimlerinin \%7,65'i artık turda olmayan kenarlara düşer;
vekilde bu oran \%0,05'tir.
Vekil saniyede toplam 2.288 vuruş yapar (GNN 1.991, kNN 1.640).
Rastgele vuruş daha hızlıdır (3.252), ama kabul oranı \%1,05'te kalır; vekilde bu oran
\%1,41'dir.
Verimlilik tek açıklama değildir: skorlaması GNN'den ucuz olan kNN (model payı \%2,83) ondan iyi
değildir; kNN'nin daha yüksek kabul oranı (\%1,88) vuruş başına yerel arama işini artırır.
% kaynak: E697, E698, E701-E703, E705-E709, E711, E712

\begin{table}[!htbp]
\centering\small\setlength{\tabcolsep}{3pt}
\caption{Mekanizma ve verimlilik, holdout2 (540 koşu/strateji). Model: skor süresinin duvar-saatine oranı; Bayat: turda artık olmayan havuz kenarına düşen kesimler. Renk: sütunun en iyi/en kötüsü.}
\label{tab:v3-mechanism}
\begin{tabular*}{\textwidth}{@{\extracolsep{\fill}}lccccc}
\toprule
Strateji & Vuruş/sn & Kabul [\%] & Model [\%] & Bayat [\%] & Gap [\%] \\
\midrule
Rastgele & 3.252 & 1{,}05 & 0{,}00 & 0{,}00 & 1{,}094 \\
kNN (0,04/0,75) & 1.640 & 1{,}88 & 2{,}83 & 0{,}18 & 1{,}015 \\
Fazla-uzunluk (0,06/0,5) & 2.288 & 1{,}41 & \iyi{0{,}65} & \iyi{0{,}05} & \iyi{0{,}922} \\
GNN (0,04/0,5) & 1.991 & 1{,}67 & \kotu{8{,}77} & \kotu{7{,}65} & 1{,}041 \\
Kazanç-MLP & 1.736 & 1{,}80 & 4{,}86 & 0{,}65 & 1{,}005 \\
Üyelik-MLP & 1.737 & 1{,}86 & 4{,}73 & 0{,}65 & 1{,}011 \\
Sabit doğrusal & 1.751 & 1{,}81 & 4{,}22 & 0{,}40 & 1{,}012 \\
Çevrimiçi (ılık) & 4.357 & 0{,}75 & 4{,}30 & 0{,}07 & \kotu{1{,}303} \\
\midrule
\multicolumn{6}{l}{\textit{Çevrimiçi ağırlık kayması (54 problem, tohum 0): ortalama / medyan / [Q1; Q3]}} \\
\quad Uzunluk: tüm tur kenarları & 0{,}623 & 0{,}609 & \multicolumn{2}{c}{[0{,}602; 0{,}653]} & \\
\quad Uzunluk: ön-bilgi havuzu & 1{,}654 & 1{,}431 & \multicolumn{2}{c}{[1{,}296; 1{,}954]} & \\
\quad Uzunluk: fazla-uz. havuzu & 1{,}614 & 1{,}388 & \multicolumn{2}{c}{[1{,}264; 1{,}918]} & \\
\quad Uzunluk: öğrenilmiş havuz & 0{,}482 & 0{,}487 & \multicolumn{2}{c}{[0{,}445; 0{,}527]} & \\
\quad Örtüşme: ön-bilgi $\cap$ fazla-uz. & 0{,}664 & 0{,}647 & \multicolumn{2}{c}{[0{,}581; 0{,}762]} & \\
\quad Örtüşme: öğrenilmiş $\cap$ fazla-uz. & 0{,}008 & 0{,}001 & \multicolumn{2}{c}{[0{,}000; 0{,}010]} & \\
\quad $\cos(w_0, w_{\text{son}})$ & $-$0{,}518 & $-$0{,}605 & \multicolumn{2}{c}{[$-$0{,}733; $-$0{,}311]} & \\
\bottomrule
\end{tabular*}
\end{table}


\paragraph{Çevrimiçi model ne öğreniyor?}
Holdout2'de güncellemeler başlangıç ağırlıklarını tersine çevirir (kosinüs ortalama $-0{,}52$,
50/54 problemde negatif) ve öğrenilmiş havuz kısa kenarları seçer: göreli uzunluk (kenar
uzunluğunun, uç noktaların 5. en yakın komşu uzaklıkları ortalamasına oranı) öğrenilmiş havuzda
$0{,}48$, ön-bilgi havuzunda $1{,}65$'tir (Şekil~\ref{fig:v3-online-drift}).
Kabul oranı sabit modelin \%1,81'ine karşı \%0,75'tir ve 54 problemin hepsinde düşüktür.
Tablo~\ref{tab:v3-online-reward} bunun nedenini geliştirme setinde ölçer (keşifsel). Öğrenici kendi
amacına ulaşır: vuruş başı ortalama ödülü 12 problemin 12'sinde sabit modelden yüksektir.
Bunu, yerel aramanın aynı maliyete onardığı kısa kenarlara yönelerek yapar: yerel arama öncesindeki
anlık hasar iki modelde aynıyken nötr vuruşların payı \%47'den \%88'e çıkar.
İyileştiren vuruşlar seyrekleşir (\%0,31'den \%0,13'e), kabul edilen bir iyileşmenin ödülü ise
değişmez.
Yani ödül (Denklem~\eqref{eq:reward}) son kaliteyle hizalı değildir. Ödül işareti birim testleriyle
doğrulandığından örüntü bir işaret hatasından kaynaklanmaz.
Geliştirmedeki son gap farkı bu 12 problemde anlamlı değildir ($p=0{,}064$): geliştirme koşuları
mekanizmayı, holdout2 (V1) zararı gösterir.
% kaynak: E599, E722, E726, E730, E732, E736, E737, E938, E941, E943, E944, E947, E948, E959

\begin{table}[!htbp]
\centering\small\setlength{\tabcolsep}{3pt}
\caption{Çevrimiçi ve sabit doğrusal skorun vuruş düzeyi ödülü (geliştirme, 12 VLSI problemi $\times$ 3 tohum, 10~sn; keşifsel). Ç$>$S: çevrimiçi değeri büyük problem sayısı.}
\label{tab:v3-online-reward}
\begin{tabular*}{\textwidth}{@{\extracolsep{\fill}}lcccc}
\toprule
Büyüklük & Çevrimiçi & Sabit & Ç$>$S & Wilcoxon $p$ \\
\midrule
Vuruş başı ödül & $-$2{,}80 & $-$25{,}88 & 12/12 & $4{,}9\times10^{-4}$ \\
Nötr vuruş [\%] & 88{,}1 & 47{,}2 & 12/12 & $4{,}9\times10^{-4}$ \\
Üç kesik kenar geri gelir [\%] & 92{,}4 & 42{,}7 & 12/12 & $4{,}9\times10^{-4}$ \\
İyileştiren vuruş [\%] & 0{,}13 & 0{,}31 & 0/12 & $4{,}9\times10^{-4}$ \\
Kötüleştiren vuruş [\%] & 11{,}8 & 52{,}5 & 0/12 & $4{,}9\times10^{-4}$ \\
Anlık hasar (yerel arama öncesi) & 194{,}8 & 194{,}3 & 8/12 & 0{,}910 \\
Ödül | iyileştiren & 1{,}90 & 1{,}93 & 5/12 & 0{,}791 \\
Havuz kenar uzunluğu / medyan & 1{,}00 & 3{,}35 & 0/12 & $4{,}9\times10^{-4}$ \\
\midrule
Son gap [\%] & 0{,}940 & 0{,}763 & 9/12 & 0{,}064 \\
Vuruş / koşu & 123.078 & 55.406 & -- & -- \\
Kabul / koşu & 145{,}1 & 145{,}0 & -- & -- \\
\bottomrule
\end{tabular*}
\end{table}


\begin{figure}[!htbp]
\centering
\begin{tikzpicture}
\begin{groupplot}[group style={group size=1 by 2, vertical sep=1.9cm},
  width=\textwidth, height=5.0cm, xmode=log, xmin=1450, xmax=10500,
  xtick={1500,2000,3000,4000,6000,8000,10000}, xticklabels={1500,2000,3000,4000,6000,8000,10000},
  xlabel={Instance size $n$ (log scale)}, tick label style={font=\scriptsize}, label style={font=\small},
  title style={font=\small, yshift=-1ex},
  ymajorgrids, grid style={black!10},
  legend style={font=\scriptsize, draw=black!40, fill=white, legend columns=2}]
\nextgroupplot[title={(a) Which edges does the kick pool target?}, ylabel={Relative edge length},
  ymin=0, ymax=3.9, legend style={at={(0.99,0.97)}, anchor=north east}]
\addplot[only marks, mark=square*, mark size=1.4pt, black!75] table[x=n, y=rel_prior] {figures/data/v3_online_drift.dat};
\addlegendentry{Prior weights $w_0$ (static)}
\addplot[only marks, mark=*, mark size=1.4pt, iyiRenk] table[x=n, y=rel_excess] {figures/data/v3_online_drift.dat};
\addlegendentry{Excess proxy}
\addplot[only marks, mark=triangle*, mark size=1.7pt, kotuRenk] table[x=n, y=rel_learned] {figures/data/v3_online_drift.dat};
\addlegendentry{Learned weights $w_{\mathrm{final}}$}
\addplot[only marks, mark=x, mark size=1.4pt, black!40] table[x=n, y=rel_all] {figures/data/v3_online_drift.dat};
\addlegendentry{All tour edges}
\nextgroupplot[title={(b) Kick acceptance rate}, ylabel={Accepted kicks [\%] (log)}, ymode=log, ymin=0.1, ymax=40,
  log ticks with fixed point, legend style={at={(0.01,0.97)}, anchor=north west}]
\addplot[only marks, mark=square*, mark size=1.4pt, black!75] table[x=n, y=acc_linprior] {figures/data/v3_online_drift.dat};
\addlegendentry{Static-linear}
\addplot[only marks, mark=*, mark size=1.4pt, iyiRenk] table[x=n, y=acc_excess] {figures/data/v3_online_drift.dat};
\addlegendentry{Excess}
\addplot[only marks, mark=triangle*, mark size=1.7pt, kotuRenk] table[x=n, y=acc_online] {figures/data/v3_online_drift.dat};
\addlegendentry{Online (warm)}
\addplot[only marks, mark=x, mark size=1.4pt, black!40] table[x=n, y=acc_random] {figures/data/v3_online_drift.dat};
\addlegendentry{Random}
\end{groupplot}
\end{tikzpicture}

\caption{Çevrimiçi (ılık) modelin havuz kayması, holdout2 (tohum 0): (a) havuzların hedeflediği kenarların göreli uzunluğu, (b) kabul edilen vuruş oranı; yatay eksen problem boyutu.}
\label{fig:v3-online-drift}
\end{figure}

\subsection{Duyarlılık: havuz kuralı ve bütçe}\label{sec:bulgular-duyarlilik}

Şekil~\ref{fig:v3-pool-grid}, geliştirme setlerinde önceden yazılmış kuralla seçilen havuz
ayarlarını (Bölüm~\ref{sec:deney-asama3}) ızgara içinde gösterir. Seçilen ayar, kNN, fazla-uzunluk
ve GNN için varsayılandan anlamlı farklı değildir ($p=0{,}82$, $0{,}36$, $0{,}53$).
Izgara $f\ge0{,}04$'te geniş bir plato gösterir. Küçük havuz ($f=0{,}02$) yalnız hedefleme olasılığı
$0{,}75$ ve üzerindeyken geride kalır; olasılık $0{,}9$'da üç sinyal de rastgeleden kötüdür
($+0{,}067$ ile $+0{,}082$ puan arası).
% kaynak: E750-E753, E788

\begin{figure}[!htbp]
\centering
\begin{tikzpicture}
\begin{groupplot}[group style={group size=3 by 1, horizontal sep=0.35cm},
  width=0.385\textwidth, height=5.0cm, xmin=0.01, xmax=0.11, ymin=-0.17, ymax=0.11,
  xtick={0.02,0.04,0.06,0.10}, xticklabels={0.02,0.04,0.06,0.10},
  xlabel={Pool fraction $f$}, tick label style={font=\scriptsize}, label style={font=\small},
  title style={font=\small, yshift=-1ex},
  y tick label style={/pgf/number format/fixed, /pgf/number format/precision=2},
  ymajorgrids, grid style={black!10}]
\nextgroupplot[title={kNN}, ylabel={$\Delta$ gap vs.\ Random [pt]}, legend to name=gridleg, legend columns=4, legend style={font=\scriptsize, draw=black!40}, ]
\addplot[black!45, dashed, forget plot] coordinates {(0.01,0) (0.11,0)};
\addplot[thick, iyiRenk, mark=*, mark size=1.4pt] table[x=frac, y=p050] {figures/data/v3_pool_grid_knn.dat};
\addlegendentry{$p=0.5$}
\addplot[thick, orange!85!black, dashed, mark=triangle*, mark size=1.6pt] table[x=frac, y=p075] {figures/data/v3_pool_grid_knn.dat};
\addlegendentry{$p=0.75$}
\addplot[thick, kotuRenk, densely dotted, mark=square*, mark size=1.3pt] table[x=frac, y=p090] {figures/data/v3_pool_grid_knn.dat};
\addlegendentry{$p=0.9$}
\addplot[only marks, mark=star, mark size=4pt, very thick, black] table[x=frac, y=d] {figures/data/v3_pool_grid_knn_sel.dat};
\addlegendentry{selected rule}
\node[anchor=south east, font=\scriptsize, fill=white, inner sep=1pt, text=black!60] at (axis cs:0.105,0) {Random};
\nextgroupplot[title={Excess}, yticklabels={}, ]
\addplot[black!45, dashed, forget plot] coordinates {(0.01,0) (0.11,0)};
\addplot[thick, iyiRenk, mark=*, mark size=1.4pt] table[x=frac, y=p050] {figures/data/v3_pool_grid_excess.dat};
\addplot[thick, orange!85!black, dashed, mark=triangle*, mark size=1.6pt] table[x=frac, y=p075] {figures/data/v3_pool_grid_excess.dat};
\addplot[thick, kotuRenk, densely dotted, mark=square*, mark size=1.3pt] table[x=frac, y=p090] {figures/data/v3_pool_grid_excess.dat};
\addplot[only marks, mark=star, mark size=4pt, very thick, black] table[x=frac, y=d] {figures/data/v3_pool_grid_excess_sel.dat};
\node[anchor=south east, font=\scriptsize, fill=white, inner sep=1pt, text=black!60] at (axis cs:0.105,0) {Random};
\nextgroupplot[title={GNN}, yticklabels={}, ]
\addplot[black!45, dashed, forget plot] coordinates {(0.01,0) (0.11,0)};
\addplot[thick, iyiRenk, mark=*, mark size=1.4pt] table[x=frac, y=p050] {figures/data/v3_pool_grid_gnn.dat};
\addplot[thick, orange!85!black, dashed, mark=triangle*, mark size=1.6pt] table[x=frac, y=p075] {figures/data/v3_pool_grid_gnn.dat};
\addplot[thick, kotuRenk, densely dotted, mark=square*, mark size=1.3pt] table[x=frac, y=p090] {figures/data/v3_pool_grid_gnn.dat};
\addplot[only marks, mark=star, mark size=4pt, very thick, black] table[x=frac, y=d] {figures/data/v3_pool_grid_gnn_sel.dat};
\node[anchor=south east, font=\scriptsize, fill=white, inner sep=1pt, text=black!60] at (axis cs:0.105,0) {Random};
\end{groupplot}
\node[anchor=north] at ([yshift=-0.9cm]group c2r1.south) {\pgfplotslegendfromname{gridleg}};
\end{tikzpicture}

\caption{Havuz kuralı ızgarası (geliştirme setleri, keşifsel): üç sinyal için havuz oranı $f$ ve hedefleme olasılığına göre rastgeleye fark; yıldız seçilen kuraldır.}
\label{fig:v3-pool-grid}
\end{figure}

Vekilin GNN'ye karşı avantajı erken dönemde büyüktür ve bütçe uzadıkça daralır: fark 0,1~sn'de
$-0{,}44$, 1~sn'de $-0{,}25$, 10~sn'de $-0{,}12$ puandır.
Rastgeleye karşı fark daha kararlıdır (aynı noktalarda $-0{,}17$, $-0{,}28$, $-0{,}17$).
Erken dönemdeki GNN farkının kaynağı model maliyetidir: GNN 0,1~sn'de rastgeleden bile
$0{,}277$ puan kötüdür ve rastgeleden ancak 5~sn'den sonra ayrışır ($-0{,}060$
puan, düzeltmesiz $p=0{,}001$).
Çevrimiçi modelin rastgele vuruşa göre zararı 1~sn'de en büyüktür ($+0{,}369$ puan) ve bütçe
uzadıkça azalarak 10~sn'de $+0{,}209$ puana iner (keşifsel, doğrulayıcı koşunun ara noktaları; V1'deki
$+0{,}291$ puan ise sabit doğrusal modele göre farktır).
% kaynak: E683, E684, E689, E690, E695, E696, E991, E1016, E1026, E1037

Yerel ölçek düzeltmesinin ablasyonu olarak, protokol sonrasında keşifsel bir ham uzunluk kolu (Bölüm~\ref{sec:yontem-sinyal}), rastgele ve
fazla-uzunluk kollarıyla aynı oturumda ve aynı ayarlarla koşuldu. Ham uzunluk rastgeleye
göre gap'i $0{,}189$ puan düşürdü (49/54, $p=1{,}5\times10^{-9}$).
Ölçek düzeltmeli vekille arasında anlamlı bir ortalama fark gözlenmedi ($+0{,}009$ puan, \%95 GA
$[-0{,}032;\,+0{,}053]$, 25'e 28, $p=0{,}99$); güven aralığının $0{,}05$ eşiğine yakınlığı nedeniyle
resmi eşdeğerlik teyidi gelecek ön-kayıtlı bir çalışmaya bırakılmıştır.
Aynı problem ve tohumlarla ayrı oturumda yeniden ölçülen fazla-uzunluk$-$rastgele farkı
$-0{,}197$ puandır (V2: $-0{,}172$).
Ölçek düzeltmesi kümeli problemlerde ortalamada vekil lehine ($0{,}089$ puan; ham kol 14 problemin
6'sında önde), üç ulusal problemde ham uzunluk lehine ($0{,}239$ puan) sonuç verdi; diğer alt
kümelerde farklar $0{,}05$ puanın altındaydı (betimsel).
% kaynak: E600, E965, E972-E979

\subsection{Keşifsel sağlamlık blokları}\label{sec:bulgular-saglamlik}

Tablo~\ref{tab:v4-robustness}, koşulardan önce kaydedilen üç keşifsel bloğun sonuçlarını verir
(Bölüm~\ref{sec:deney-saglamlik}).

\paragraph{Klasik yerel vuruş.}
Hedefli global vuruş, rastgele segment vuruşundan ortalama $1{,}038$ puan düşük gap verir (51/54,
$p_{\mathrm{Holm}}=1{,}3\times10^{-9}$). Fark çarpıktır: en büyük katkı kümeli problemlerden gelir
($-3{,}10$ puan), Hodges--Lehmann kestirimi $-0{,}535$ puandır.
Bu motorda ve bütçede segment geometrisi, rastgele kesimlerle de globalden kötüdür ($+0{,}784$ puan,
7/54); yalnız üç ulusal problemde yerel geometri öndedir.
Segment vuruşu saniyede yaklaşık iki buçuk kat fazla vuruş yapar (5.898'e karşı 2.225), ama kabul
oranı \%0,55'tir (global \%1,41).
Yerel geometride hedeflemenin ölçülebilir etkisi yoktur (fazla-uzunluk $+0{,}037$, ham uzunluk
$-0{,}014$ puan; ikisinde de $p_{\mathrm{Holm}}=1{,}0$). Burada havuzdan yalnız ilk kesim seçilir,
diğer ikisi pencere içinde rastgeledir.
Doğrulama setine veri sızıntısını önlemek amacıyla pencere genişliği ($W=50$) sabit tutulmuştur;
elde edilen sonuç, segment vuruşunun LK tabanlı çözücülerdeki yerleşik performansından
\citep{applegate_2003_clk} bağımsız olarak, 2-opt/Or-opt motorundaki göreli davranışını karakterize eder.
% kaynak: E1047-E1049, E1051, E1057, E1059

\paragraph{Uzun bütçe.}
60~sn'de (5 tohum) vekil hem rastgeleden ($-0{,}071$ puan, 40/54, $p_{\mathrm{Holm}}=7{,}7\times10^{-4}$)
hem GNN'den ($-0{,}078$ puan, 35/54, $p_{\mathrm{Holm}}=0{,}030$) düşük gap verir.
Avantaj bütçe uzadıkça daralır: rastgeleye karşı 10, 30 ve 60~sn'de $-0{,}242$, $-0{,}108$ ve
$-0{,}071$; GNN'ye karşı $-0{,}149$, $-0{,}092$ ve $-0{,}078$ puandır.
GNN'ye karşı farkın \%95 GA'sı 60~sn'de de sıfırın altındadır ($[-0{,}172;\,-0{,}008]$); V3'ün
0,05 puanlık kötü-değil sonucu bu bütçede de geçerlidir.
Mutlak gap'ler de küçüldüğü için göreli kazanç daha yavaş azalır: vekil, rastgelenin gap'ini 10~sn'de
\%21 (\%1,139'dan \%0,897'ye), 60~sn'de \%11 (\%0,649'dan \%0,578'e) düşürür.
% kaynak: E1062-E1065, E1072, E1073, E1077-E1080

\paragraph{Havuz kuralı.}
Geliştirmede seçilen kuralın dört komşusundan üçü holdout2'de de rastgeleden açıkça iyidir
($-0{,}178$ ile $-0{,}221$ puan arası, 49--50/54).
Tek istisna, küçük havuzun her kesimde hedeflendiği birleşimdir ($f=0{,}03$, $p=1$): bu ayar
54 problemin yalnız 11'inde rastgeleden iyidir ve seçilen kuraldan $0{,}313$ puan geridedir.
Seçilen kural, komşularından ya iyidir ya da onlardan ayrışmaz ($+0{,}034$ ile $+0{,}076$ puan;
$f=0{,}03$, $p=0{,}25$ ile fark anlamlı değil).
Bu örüntü geliştirme ızgarasıyla aynıdır (Şekil~\ref{fig:v3-pool-grid}): geniş bir plato ve tek
zararlı bölge olarak küçük havuz ile çok yüksek hedefleme olasılığının birleşimi.
% kaynak: E1086-E1093

\begin{table}[!htbp]
\centering\small\setlength{\tabcolsep}{3pt}
\caption{Keşifsel sağlamlık blokları (holdout2; protokol koşulardan önce kaydedildi). Fark [puan]: negatif değer ilk terim lehine. Önde: ilk terimin önde olduğu problem sayısı. $p_{\mathrm{Holm}}$: blok içinde (b'de her zaman noktasında) Holm düzeltmeli Wilcoxon. Renk: $p_{\mathrm{Holm}}<0{,}05$ ve GA sıfırı dışlar.}
\label{tab:v4-robustness}
\begin{tabular*}{\textwidth}{@{\extracolsep{\fill}}lcccc}
\toprule
Karşılaştırma & Fark & \%95 GA & Önde & $p_{\mathrm{Holm}}$ \\
\midrule
\multicolumn{5}{@{}l}{\textit{(a) Klasik yerel (segment) vuruş --- 10~sn, 10 tohum}} \\
Fazla-uzunluk (global) $-$ rastgele yerel & \iyi{$-$1{,}038} & [$-$1{,}447; $-$0{,}688] & 51/54 & $1{,}3\times10^{-9}$ \\
Fazla-uzunluk (yerel) $-$ rastgele yerel & +0{,}037 & [$-$0{,}028; +0{,}117] & 32/54 & 1{,}000 \\
Ham uzunluk (yerel) $-$ rastgele yerel & $-$0{,}014 & [$-$0{,}079; +0{,}049] & 28/54 & 1{,}000 \\
Fazla-uzunluk: yerel $-$ global & \kotu{+1{,}074} & [+0{,}707; +1{,}505] & 2/54 & $2{,}9\times10^{-9}$ \\
Rastgele: yerel $-$ global & \kotu{+0{,}784} & [+0{,}462; +1{,}155] & 7/54 & $3{,}5\times10^{-6}$ \\
\midrule
\multicolumn{5}{@{}l}{\textit{(b) Uzun bütçe --- 60~sn, 5 tohum (aynı koşuların ara noktaları)}} \\
Fazla-uzunluk $-$ rastgele, 10~sn & \iyi{$-$0{,}242} & [$-$0{,}337; $-$0{,}154] & 50/54 & $1{,}3\times10^{-8}$ \\
Fazla-uzunluk $-$ GNN, 10~sn & \iyi{$-$0{,}149} & [$-$0{,}249; $-$0{,}067] & 41/54 & $7{,}6\times10^{-5}$ \\
Fazla-uzunluk $-$ rastgele, 30~sn & \iyi{$-$0{,}108} & [$-$0{,}165; $-$0{,}051] & 41/54 & $1{,}4\times10^{-5}$ \\
Fazla-uzunluk $-$ GNN, 30~sn & \iyi{$-$0{,}092} & [$-$0{,}185; $-$0{,}022] & 37/54 & 0{,}008 \\
Fazla-uzunluk $-$ rastgele, 60~sn & \iyi{$-$0{,}071} & [$-$0{,}116; $-$0{,}024] & 40/54 & $7{,}7\times10^{-4}$ \\
Fazla-uzunluk $-$ GNN, 60~sn & \iyi{$-$0{,}078} & [$-$0{,}172; $-$0{,}008] & 35/54 & 0{,}030 \\
\midrule
\multicolumn{5}{@{}l}{\textit{(c) Havuz kuralı holdout2'de --- 10~sn, 10 tohum}} \\
$f=0{,}03$, $p=0{,}25$ $-$ rastgele & \iyi{$-$0{,}221} & [$-$0{,}291; $-$0{,}157] & 50/54 & $1{,}3\times10^{-8}$ \\
$f=0{,}03$, $p=1{,}00$ $-$ rastgele & +0{,}059 & [$-$0{,}038; +0{,}146] & 11/54 & 0{,}007 \\
$f=0{,}12$, $p=0{,}25$ $-$ rastgele & \iyi{$-$0{,}178} & [$-$0{,}236; $-$0{,}127] & 50/54 & $8{,}9\times10^{-9}$ \\
$f=0{,}12$, $p=1{,}00$ $-$ rastgele & \iyi{$-$0{,}202} & [$-$0{,}278; $-$0{,}136] & 49/54 & $3{,}2\times10^{-8}$ \\
$f=0{,}03$, $p=0{,}25$ $-$ seçilen kural & +0{,}034 & [$-$0{,}001; +0{,}074] & 21/54 & 0{,}159 \\
$f=0{,}03$, $p=1{,}00$ $-$ seçilen kural & \kotu{+0{,}313} & [+0{,}262; +0{,}360] & 2/54 & $2{,}2\times10^{-8}$ \\
$f=0{,}12$, $p=0{,}25$ $-$ seçilen kural & \kotu{+0{,}076} & [+0{,}036; +0{,}122] & 19/54 & 0{,}005 \\
$f=0{,}12$, $p=1{,}00$ $-$ seçilen kural & \kotu{+0{,}052} & [+0{,}027; +0{,}079] & 14/54 & $4{,}7\times10^{-4}$ \\
\bottomrule
\end{tabular*}
\end{table}


\subsection{Hakem odaklı keşifsel bloklar}\label{sec:bulgular-hakem}

Tablo~\ref{tab:v5-reviewer}, koşulardan önce kaydedilen beş oturumun sonuçlarını verir
(Bölüm~\ref{sec:deney-hakem}).

\paragraph{Yayımlanmış öğrenilmiş skor.}
NeuroLKH'nin ileri geçişi problem başına 0,47~sn ($n=1.689$) ile 2,64~sn ($n=9.882$) arasında
sürdü. Bu süre bütçeden düşüldüğünde vekil ortalamada $0{,}098$ puan öndedir, ancak fark problemler
arasında tutarlı değildir (vekil 20/54 problemde önde, $p_{\mathrm{Holm}}=0{,}58$, medyan fark
NeuroLKH lehine $0{,}033$). Ortalamadaki farkı büyük ölçüde üç büyük ulusal problem yaratır: orada
ağın süresi bütçenin \%20--26'sını alır ve NeuroLKH rastgele vuruşun da gerisine düşer.
Süre bütçeden düşülmediğinde ortalamalar eşittir ($+0{,}009$ puan, \%95 GA $[-0{,}067;\,+0{,}069]$),
fakat NeuroLKH problemlerin çoğunda biraz öndedir (37/54, medyan $0{,}066$ puan,
$p_{\mathrm{Holm}}=0{,}034$). Ortalamayı dengeleyen, vekilin üç ulusal problemdeki üstünlüğüdür
($-0{,}84$ puan). Yani yayımlanmış ve büyük veriyle eğitilmiş bir modelin sinyali, maliyeti
yok sayıldığında, vekilden küçük ama tutarlı bir fark taşır; maliyet sayıldığında bu fark kaybolur.
Ayar eşitliği vekilin diğer rakiplere karşı konumunu değiştirmez. Ayarlı kazanç-MLP varsayılan
kuraldan ayrışmaz ($-0{,}013$ puan, $p_{\mathrm{Holm}}=0{,}58$) ve vekil ayarlı kazanç-MLP'den
($-0{,}070$ puan, 35/54, $p_{\mathrm{Holm}}=0{,}026$), GNN'den ($-0{,}119$ puan, 39/54) ve
rastgeleden ($-0{,}150$ puan, 41/54) düşük gap verir.
% kaynak: E1106-E1113, E1120, E1122

\paragraph{Eşit vuruş sayısı.}
Her kol 20.000 vuruş yaptığında ve model süresi sayılmadığında hedeflemenin değeri açıkça görülür:
bütün sinyaller rastgeleden $0{,}22$--$0{,}30$ puan iyidir ve 54 problemin 49--51'inde öndedir.
Vekilin öğrenilmiş skorlara karşı eşit süredeki üstünlüğü ise bu koşulda kaybolur. Vekil GNN ile
aynı düzeydedir ($-0{,}004$ puan, 25/54); kazanç-MLP ($+0{,}074$ puan) ve NeuroLKH ($+0{,}039$ puan)
vekilden biraz iyidir (32/54 ve 40/54), ancak ikisi de Holm düzeltmesinden sonra anlamlı değildir
($p_{\mathrm{Holm}}=0{,}075$ ve $0{,}051$). Vuruş başına en iyi sinyal eğitimsiz kNN sırasıdır: vekilden
$0{,}088$ puan iyidir (37/54, $p_{\mathrm{Holm}}=0{,}011$).
Vuruş hızları farkın kaynağını gösterir: aynı oturumda vekil saniyede 1.196 vuruş yaparken GNN 524
vuruş yapar ve sürenin \%52'sini modele harcar; kazanç-MLP 919, kNN 867, NeuroLKH 1.073 vuruş/sn'dir.
Eşit süredeki sıralama, büyük ölçüde sinyalin niteliğinden değil maliyetinden doğar.
% kaynak: E1128-E1135, E1137, E1139, E1141, E1143, E1145

\paragraph{Çevrimiçi öğrenmede kredi tanımı.}
Çevrimiçi modelin zararı kredi tanımından gelir. İşaretli ödülle model sabit ön-bilgiden $0{,}335$
puan kötüdür (2/54); ödül $\max(0,r)$ ya da $\mathbb{I}[\text{kabul}]$ olduğunda bu zarar tamamen
kaybolur. İki tanım da işaretli ödülden 54 problemin 54'ünde iyidir ($-0{,}389$ ve $-0{,}338$ puan) ve
sabit ön-bilgiden ayrışmaz ($-0{,}054$ puan, $p_{\mathrm{Holm}}=0{,}13$; $-0{,}003$ puan,
$p_{\mathrm{Holm}}=0{,}62$). İşaretli ödüllü model saniyede 2.745 vuruş yapıp yalnız \%0,8'ini kabul
ettirir; diğer iki tanımda bu değerler yaklaşık 1.560 vuruş/sn ve \%1,4'tür. Bu, modelin ucuz ve
nötr vuruşlara kaydığı mekanizmayla (Bölüm~\ref{sec:bulgular-mekanizma}) tutarlıdır. Doğru krediyle
çevrimiçi öğrenme zarar vermez, ama bu bütçede sabit modele bir şey de katmaz.
% kaynak: E1148-E1153, E1157, E1159, E1161

\paragraph{Tek işçi ve kol sırası.}
Tek işçiyle ve rastgele kol sırasıyla vekil yine rastgeleden ($-0{,}093$ puan, \%95 GA
$[-0{,}182;\,-0{,}012]$) ve GNN'den ($-0{,}057$ puan) düşük gap verir. Yalnız 2 tohum koşulduğu için iki
fark da Holm düzeltmesinden sonra anlamlı değildir. Tek işçide vuruş hızı dört işçili oturumların
yaklaşık 1,7 katıdır (rastgele 3.798'e karşı 2.180 vuruş/sn); dört işçi muhtemelen bellek ve
önbellek paylaşımı yüzünden her kolu birlikte yavaşlatır. Bir partide ilk koşan kol diğerlerinden yaklaşık
\%18 yavaş vurur (Spearman $\rho=0{,}12$). Ana koşuda sıra tohumlar arasında döndürüldüğü için bu
etki büyük ölçüde dengelenmiştir.
% kaynak: E1164, E1165, E1167, E1171, E1127, E1181

\paragraph{LKH-3.}
LKH-3'ün ILS kipinde hedefli vuruş rastgele vuruştan ayrışmaz ($+0{,}001$ puan, \%95 GA
$[-0{,}075;\,+0{,}067]$, 30/54, $p_{\mathrm{Holm}}=0{,}59$). Hedefli kol aynı sürede \%22 daha az deneme
yapar (699'a karşı 897); bunun olası nedeni, havuzun her vuruşta $O(n)$ sürede yeniden kurulmasıdır. Ayrıca ILS kipinin
kendisi LKH'nin varsayılan ayarlarından kötüdür ($+0{,}062$ puan, 10/54). LKH-3, 2-opt/Or-opt
motorundan çok daha düşük gap'lere ulaşır (\%0,26--0,32'ye karşı \%1,07--1,22). Bu düzeyde bir
yerel aramada uzun kenar bilgisi vuruş seçimine ek bir kazanç getirmez.
% kaynak: E1172-E1180

% OTOMATIK URETIM -- evidence/scripts/build_v5.py (elle duzenleme)
\begin{table}[!htbp]
\centering\small\setlength{\tabcolsep}{3pt}
\caption{Hakem odaklı keşifsel bloklar (holdout2; protokol koşulardan önce kaydedildi). Fark [puan]: negatif değer ilk terim lehine. Önde: ilk terimin önde olduğu problem sayısı. $p_{\mathrm{Holm}}$: blok içinde Holm düzeltmeli Wilcoxon. Renk: $p_{\mathrm{Holm}}<0{,}05$ ve GA sıfırı dışlar.}
\label{tab:v5-reviewer}
\begin{tabular*}{\textwidth}{@{\extracolsep{\fill}}lcccc}
\toprule
Karşılaştırma & Fark & \%95 GA & Önde & $p_{\mathrm{Holm}}$ \\
\midrule
\multicolumn{5}{@{}l}{\textit{(a) Eşit süre --- 10~sn, 10 tohum; ayarlı $(f,p)$: kazanç-MLP (0{,}04; 0{,}50), NeuroLKH (0{,}10; 0{,}75)}} \\
Fazla-uzunluk $-$ NeuroLKH (süre sayılır) & $-$0{,}098 & [$-$0{,}210; $-$0{,}008] & 20/54 & 0{,}583 \\
Fazla-uzunluk $-$ NeuroLKH (süre sayılmaz) & +0{,}009 & [$-$0{,}067; +0{,}069] & 17/54 & 0{,}034 \\
NeuroLKH (süre sayılır) $-$ rastgele & $-$0{,}052 & [$-$0{,}117; +0{,}022] & 41/54 & 0{,}002 \\
NeuroLKH (süre sayılmaz) $-$ rastgele & \iyi{$-$0{,}158} & [$-$0{,}214; $-$0{,}107] & 47/54 & $7{,}7\times10^{-7}$ \\
Fazla-uzunluk $-$ kazanç-MLP (ayarlı) & \iyi{$-$0{,}070} & [$-$0{,}128; $-$0{,}018] & 35/54 & 0{,}026 \\
Kazanç-MLP (ayarlı) $-$ kazanç-MLP (varsayılan) & $-$0{,}013 & [$-$0{,}058; +0{,}032] & 31/54 & 0{,}583 \\
Fazla-uzunluk $-$ GNN & \iyi{$-$0{,}119} & [$-$0{,}188; $-$0{,}061] & 39/54 & $8{,}6\times10^{-4}$ \\
Fazla-uzunluk $-$ rastgele & \iyi{$-$0{,}150} & [$-$0{,}218; $-$0{,}089] & 41/54 & $1{,}4\times10^{-5}$ \\
\midrule
\multicolumn{5}{@{}l}{\textit{(b) Eşit vuruş sayısı --- 20.000 vuruş, 5 tohum, model süresi sayılmaz}} \\
Fazla-uzunluk $-$ GNN & $-$0{,}004 & [$-$0{,}076; +0{,}065] & 25/54 & 0{,}633 \\
Fazla-uzunluk $-$ kazanç-MLP (ayarlı) & +0{,}074 & [+0{,}002; +0{,}166] & 22/54 & 0{,}075 \\
Fazla-uzunluk $-$ NeuroLKH & +0{,}039 & [$-$0{,}028; +0{,}106] & 14/54 & 0{,}051 \\
Fazla-uzunluk $-$ kNN sırası & \kotu{+0{,}088} & [+0{,}033; +0{,}156] & 17/54 & 0{,}011 \\
GNN $-$ rastgele & \iyi{$-$0{,}220} & [$-$0{,}306; $-$0{,}141] & 51/54 & $1{,}9\times10^{-7}$ \\
Kazanç-MLP (ayarlı) $-$ rastgele & \iyi{$-$0{,}298} & [$-$0{,}423; $-$0{,}185] & 50/54 & $3{,}8\times10^{-8}$ \\
NeuroLKH $-$ rastgele & \iyi{$-$0{,}262} & [$-$0{,}345; $-$0{,}189] & 49/54 & $1{,}1\times10^{-7}$ \\
Fazla-uzunluk $-$ rastgele & \iyi{$-$0{,}224} & [$-$0{,}294; $-$0{,}157] & 51/54 & $2{,}1\times10^{-8}$ \\
\midrule
\multicolumn{5}{@{}l}{\textit{(c) Çevrimiçi öğrenmede kredi tanımı --- 10~sn, 10 tohum}} \\
Çevrimiçi, $\max(0,r)$ $-$ doğrusal ön-bilgi & $-$0{,}054 & [$-$0{,}111; $-$0{,}004] & 31/54 & 0{,}130 \\
Çevrimiçi, $\mathbb{I}[\text{kabul}]$ $-$ doğrusal ön-bilgi & $-$0{,}003 & [$-$0{,}072; +0{,}062] & 29/54 & 0{,}621 \\
Çevrimiçi, işaretli ödül $-$ doğrusal ön-bilgi & \kotu{+0{,}335} & [+0{,}245; +0{,}432] & 2/54 & $9{,}8\times10^{-10}$ \\
Çevrimiçi, $\max(0,r)$ $-$ çevrimiçi, işaretli ödül & \iyi{$-$0{,}389} & [$-$0{,}494; $-$0{,}292] & 54/54 & $9{,}8\times10^{-10}$ \\
Çevrimiçi, $\mathbb{I}[\text{kabul}]$ $-$ çevrimiçi, işaretli ödül & \iyi{$-$0{,}338} & [$-$0{,}423; $-$0{,}261] & 54/54 & $9{,}8\times10^{-10}$ \\
Doğrusal ön-bilgi $-$ rastgele & $-$0{,}081 & [$-$0{,}146; $-$0{,}023] & 32/54 & 0{,}106 \\
\midrule
\multicolumn{5}{@{}l}{\textit{(d) Tek işçi, rastgele kol sırası --- 10~sn, 2 tohum}} \\
Fazla-uzunluk $-$ rastgele & $-$0{,}093 & [$-$0{,}182; $-$0{,}012] & 29/54 & 0{,}195 \\
Fazla-uzunluk $-$ GNN & $-$0{,}057 & [$-$0{,}128; +0{,}007] & 30/54 & 0{,}355 \\
\midrule
\multicolumn{5}{@{}l}{\textit{(e) LKH-3.0.13 içinde vuruş hedefleme --- 10~sn, 5 tohum}} \\
LKH-3 ILS, hedefli vuruş $-$ LKH-3 ILS, rastgele vuruş & +0{,}001 & [$-$0{,}075; +0{,}067] & 30/54 & 0{,}592 \\
LKH-3 ILS, rastgele vuruş $-$ LKH-3 varsayılan & \kotu{+0{,}062} & [+0{,}018; +0{,}110] & 10/54 & $2{,}9\times10^{-5}$ \\
LKH-3 ILS, hedefli vuruş $-$ LKH-3 varsayılan & +0{,}063 & [$-$0{,}009; +0{,}137] & 15/54 & 0{,}001 \\
\bottomrule
\end{tabular*}
\end{table}


\subsection{Aşama 2: kestirim doğruluğu arama faydası değildir}\label{sec:bulgular-asama2}

Aşama~2'de (38 VLSI problemi, 10~sn) GNN hedeflemesi rastgeleye göre küçük ama anlamlı bir yarar
sağladı ($-0{,}052$ puan, 29/38, $p_{\mathrm{Holm}}=0{,}0004$; Tablo~\ref{tab:v2-primary-tests}).
Eğitimsiz kNN sırası ise GNN'den iyiydi (GNN$-$kNN $+0{,}049$ puan, kNN 29/38 önde,
$p_{\mathrm{Holm}}=0{,}0039$).
Bu sıralama kestirim doğruluğuyla çelişir: GNN, LKH-3 (60~sn) referans turunda olmayan kenarları kNN'den
daha iyi ayırt eder (alıcı işletim karakteristiği eğrisinin altındaki alan, AUC, $0{,}854$'e karşı $0{,}821$; 58 problemin 57'sinde;
Tablo~\ref{tab:v2-mechanism}).
Havuz bileşimi bu çelişkiyle uyumludur. Üyelik için eğitilen skorlar havuza, düzeltildiğinde az
kazandıran incelikli hataları alır; kNN havuzundaki yanlış kenarlar daha pahalıdır (referans tura
göre fazla uzunluk kNN'de $1{,}65$, GNN'de $1{,}30$; kNN 57 problemin 57'sinde önde).
Aynı MLP'yi kazanç hedefiyle eğitmek, kNN karşısındaki kaybın yaklaşık dörtte üçünü kapattı
(keşifsel, deney sonuçlar görüldükten sonra tasarlandı; kalan fark $+0{,}016$ puan, $p=0{,}35$).
Holdout2 bu tabloyu kısmen değiştirir: kazanç-MLP ile üyelik-MLP ayrışmaz ($-0{,}006$ puan,
$p_{\mathrm{Holm}}=1{,}0$), Aşama~2'de rastgeleden iyi olan kNN ($-0{,}101$ puan, 34/38) ise
holdout2'de anlamlılığını yitirir.
Dolayısıyla holdout2 bulguları, hedef uyumu, düşük hesaplama maliyeti ve havuz tazeliğinin tekil
etkenlerden ziyade birlikte çalışan tamamlayıcı mekanizmalar olarak değerlendirilmesi gerektiğini
göstermektedir.
% kaynak: E603, E606, E797, E803, E805, E807, E889, E890, E922, E924, E926, E928, E930

% OTOMATIK URETIM -- evidence/scripts/build_v2.py
\begin{table}[!htbp]
\centering\small\setlength{\tabcolsep}{3pt}
\caption{Birincil sette ön-kayıtlı hipotezler (H1--H4, tek Holm ailesi) ve ikincil karşılaştırmalar. $\Delta$: A$-$B gap farkı [puan], negatif A lehine; TOST sınırı $\pm0{,}10$.}
\label{tab:v2-primary-tests}
\begin{tabular*}{\textwidth}{@{\extracolsep{\fill}}lcccccc}
\toprule
Karşılaştırma (A$-$B) & $\Delta$ & \%95 GA & A/E/B & Wilcoxon $p$ & Holm-$p$ & TOST $p$ \\
\midrule
GNN$-$Rastgele (global) & $-$0{,}052 & [$-$0{,}077; $-$0{,}029] & 29/0/9 & $<10^{-4}$ & 0{,}0004 & 0{,}0002 \\
GNN$-$Rastgele (yerel) & +0{,}019 & [$-$0{,}010; +0{,}047] & 15/0/23 & 0{,}1657 & 0{,}3053 & $<10^{-4}$ \\
GNN$-$kNN (global) & +0{,}049 & [+0{,}020; +0{,}077] & 9/0/29 & 0{,}0013 & 0{,}0039 & 0{,}0007 \\
GNN$-$kNN (yerel) & +0{,}014 & [$-$0{,}008; +0{,}035] & 14/0/24 & 0{,}1526 & 0{,}3053 & $<10^{-4}$ \\
\midrule
kNN$-$Rastgele (global) & $-$0{,}101 & [$-$0{,}126; $-$0{,}077] & 34/0/4 & $<10^{-4}$ & --- & 0{,}5349 \\
kNN$-$Rastgele (yerel) & +0{,}004 & [$-$0{,}025; +0{,}034] & 20/0/18 & 0{,}8522 & --- & $<10^{-4}$ \\
Yerel$-$global (rastgele) & +0{,}362 & [+0{,}321; +0{,}402] & 1/0/37 & $<10^{-4}$ & --- & 1{,}0000 \\
\bottomrule
\end{tabular*}
\end{table}

\begin{table}[!htbp]
\centering\small\setlength{\tabcolsep}{2.5pt}
\caption{LKH-3 referans turunda bulunmayan (``yanlış'') tur kenarlarını ayırt etme, VLSI ortalamaları. Taban: yanlış kenar oranı; Kesinlik: havuzdaki yanlış oranı; $\rho$: GNN--kNN Spearman.}
\label{tab:v2-mechanism}
\begin{tabular*}{\textwidth}{@{\extracolsep{\fill}}lccccccccc}
\toprule
 & Taban & \multicolumn{4}{c}{AUC} & \multicolumn{2}{c}{Kesinlik [\%]} & & Örtüşme \\
\cmidrule(lr){3-6}\cmidrule(lr){7-8}
Tur durumu & [\%] & GNN & kNN & Göreli uz. & Uzunluk & GNN & kNN & $\rho$ & [\%] \\
\midrule
Başlangıç turu (58) & 20{,}5 & 0{,}878 & 0{,}852 & 0{,}816 & 0{,}789 & 85{,}2 & 82{,}1 & 0{,}75 & 65 \\
2 sn ILS turu (58) & 15{,}3 & 0{,}854 & 0{,}821 & 0{,}784 & 0{,}751 & 64{,}6 & 60{,}7 & 0{,}80 & 62 \\
\bottomrule
\end{tabular*}
\end{table}


\section{Tartışma}\label{sec:tartisma}

\paragraph{Ucuz bir vekil neden önde kalır?}
Gözlenen performans profili, birbirini destekleyen üç temel mekanizmayla tutarlıdır. Birincisi
hedef uyumu hipotezidir. Double-bridge vuruşunun getirisi, kırılan kenarın yerine daha kısa
bağlantılar kurulabilmesine bağlıdır; uzun bir kenar, özellikle uç noktalarının komşu ölçeğine
göre uzunsa, tam da bu kazancı vaat eder. Üyelik olasılığı ise ``bu kenar optimal turda mı?''
sorusunu yanıtlar ve havuzu, düzeltildiğinde az kazandıran hatalarla doldurabilir. Aşama~2 bu
hipotezi desteklerken holdout2'de GNN ile kNN ve kazanç-MLP ile üyelik-MLP ayrışmamıştır
(ikisinde de $p_{\mathrm{Holm}}=1{,}0$; Bölüm~\ref{sec:bulgular-asama2}).
İkincisi maliyettir: skor sabit zamanlıdır ve bütçenin neredeyse tamamı vuruşa ve yerel aramaya
kalır. Üçüncüsü tazeliktir: ucuz skor her tur değişiminde yeniden hesaplanabildiği için kesimler
güncel tura düşer. Bu yaklaşım, literatürdeki köklü sezgisellerle kuramsal bir süreklilik
içindedir. GLS ceza hedefini kenar maliyetine göre seçerek uzun kenarları önceliklendirir
\citep{voudouris_1999_gls}; rotalamada istenmeyen kenar özniteliklerinden türetilen bilginin
pertürbasyonu yönlendirebildiği de gösterilmiştir \citep{arnold_2019_kgls}. Desen, basit tabanların
öğrenilmiş rehberlerle rekabet ettiği önceki gözlemlerle aynı yöndedir
\citep{xia_2024_rethinking, bother_2022_wrong}.
% kaynak: E605, E606
Eşit vuruş oturumu bu üç mekanizmayı birbirinden ayırır (Bölüm~\ref{sec:bulgular-hakem}). Vuruş
sayısı eşitlendiğinde öğrenilmiş skorlar vekilden kötü değildir; kazanç-MLP ve NeuroLKH biraz öndedir.
Eşit süredeki sıralama bu yüzden ağırlıkla maliyet ve tazelikten doğar; hedef uyumu tek başına
vekilin üstünlüğünü açıklamaz.
% kaynak: E1128-E1130, E1137, E1139

\paragraph{Hangi eğitimsiz sinyal?}
Eğitimsiz olmak tek başına yetmez. kNN sırası holdout2'de rastgeleden anlamlı ayrışmadı ve karışık
problemlerde gerisinde kaldı (betimsel); fazla-uzunluk ise kNN'yi 54 problemin 43'ünde geçti
(keşifsel).
kNN bir kenarın komşuluk \emph{sırasını}, fazla-uzunluk ise kenarın \emph{büyüklüğünü} ölçer;
her iki kol geliştirme setinde kendi en iyi havuz kuralıyla yapılandırıldığından, bu karşılaştırma
pratik uygulama düzeyindeki uçtan uca strateji başarımını yansıtır.
Keşifsel ham uzunluk kolu, yerel ölçek düzeltmesinin ablasyonudur ve ``önce en basit vekile karşı sına'' ilkesini kendi vekilimize uygular;
büyüklük bilgisinin temel belirleyici olduğunu düşündürür: yerel ölçek düzeltmesi olmadan en uzun
kenarları hedeflemek de rastgeleye karşı benzer kazanç sağladı ve vekilden ortalamada ayrışmadı
(fark $+0{,}009$ puan, \%95 GA $[-0{,}032;\,+0{,}053]$).
Eşit vuruş oturumu bu tabloyu inceltir: vuruş başına en iyi sinyal kNN sırasıdır (vekilden
$0{,}088$ puan iyi, 37/54). kNN'nin eşit sürede geride kalması, sinyalinden çok vuruş hızından
kaynaklanır (aynı oturumda 867'ye karşı 1.196 vuruş/sn). İki eğitimsiz sinyal arasındaki seçim bu
nedenle bir nitelik--maliyet takasıdır.
Uygulayıcı için sonuç sadedir: bu motorda ve 10~sn'de, komşu listesi dahi gerektirmeyen ham kenar
uzunlukları güçlü bir başlangıç noktasıdır. Yerel ölçek düzeltmesi ortalamada ek kazanç üretmemiş
olup, kümeli yapılardaki olası avantajı bağımsız bir doğrulayıcı hipotez olarak değerlendirilebilir.
% kaynak: E604, E750, E751, E965, E972, E973, E976, E1131, E1137, E1145

\paragraph{Hedefleme nerede işe yarar?}
Kazanç kümeli ve gerçek problemlerde toplanır, uniform problemlerde küçük kalır: hedefleme bir
``kötü kenar'' sinyaline dayanır ve bu sinyal ancak kenar uzunlukları heterojen olduğunda bilgi
taşır. Katman etiketinden bağımsız kümelenme göstergesi de bunu destekler: kazanç, en yakın komşu
uzaklığının değişim katsayısıyla birlikte büyür ($\rho=-0{,}56$).
Yalnız uniform örneklerden oluşan bir kıyas kümesi bu yüzden hedeflemenin değerini
olduğundan küçük gösterir.
% kaynak: E1038

\paragraph{Yayımlanmış bir modele ve daha güçlü bir yerel aramaya karşı.}
NeuroLKH'nin sinyali, süresi sayılmadığında problemlerin çoğunda vekilden küçük bir farkla iyidir
(medyan $0{,}066$ puan). Öğrenilmiş bir modelin taşıdığı ek bilgi gerçektir, ama küçüktür; tek
seferlik ve kısa bütçeli bir çözümde ağın süresi bu kazancı siler. Model çıktısının birçok koşuya
yayılabildiği durumlarda, örneğin aynı problemin tekrar tekrar çözülmesinde, NeuroLKH gibi bir
model küçük bir kazanç sağlayabilir. Tek seferlik kullanımda eğitimsiz vekil daha basittir ve
ondan istatistiksel olarak ayrışmaz.
LKH-3'te ise hedefleme ölçülebilir bir kazanç sağlamadı. Kayıtta öngörülen olası açıklama şudur:
LK'nin derin hamleleri uzun kenarların çoğunu zaten kendisi onarır. Vuruş hedeflemenin değeri bu
yüzden yerel aramanın gücüne bağlıdır; bulgular hızlı ve görece zayıf yerel aramalı ILS için
geçerlidir.
% kaynak: E1106, E1107, E1172, E1180

\paragraph{Çevrimiçi öğrenme ve ödül tasarımı.}
Çevrimiçi model başarısız bir öğrenici değildir; kendisine verilen amacı yükseltir. Sorun amaçtadır:
vuruşların çoğu turu iyileştirmediğinde vuruş başına ortalama ödül, yerel aramanın onardığı nötr
vuruşu ödüllendirir (Bölüm~\ref{sec:bulgular-mekanizma}). Uyarlamalı operatör seçimi literatürü, kredi atamasının seçim başarımını etkilediğini \citep{fialho_2010_aos} ve nadir ama büyük iyileşmeleri öne çıkaran uç-değer kredisinin ortalama krediye göre avantajlı olabileceğini \citep{fialho_2008_extreme} öne sürer. Geliştirme setindeki ölçümümüz (12 problem, keşifsel) bunun bir örneğidir. Hakem odaklı blok bunu
holdout2'de doğrudan sınadı: kredi $\max(0,r)$ ya da kabul göstergesi olduğunda zarar 54 problemin
54'ünde kaybolur ve model sabit ön-bilgi düzeyine döner. Tasarım dersi şudur: seyrek iyileşmeli
aramalarda kredi, işaretli ortalama maliyet değişimi yerine kabul edilen iyileşmelere dayanmalı ve
öğrenicinin kendi ödülü son kaliteyle birlikte izlenmelidir. Zarar öğrenmenin kendisinden değil
kredi tanımından gelir; doğru krediyle de bu bütçede sabit modele bir kazanç eklenmedi.
% kaynak: E1148-E1152

\paragraph{Uygulayıcı için öneriler.}
\begin{itemize}
\item Vuruşları en uzun tur kenarlarına yönlendirin (ham ya da uç noktaların yakın-komşu ölçeğine
göre düzeltilmiş uzunluk); skor sabit zamanlıdır, eğitim ve çözücü etiketi gerektirmez. Bu tür
motorlarda hedefli global double-bridge, rastgele segment vuruşundan belirgin biçimde iyidir.
% kaynak: E1047
\item Havuzu tur kenarlarının \%4--6'sı kadar, hedefleme olasılığını $0{,}5$--$0{,}75$ aralığında
tutun; yalnız küçük havuz ile çok yüksek olasılığın birleşiminden kaçının. Bu plato holdout2'de de
görüldü.
% kaynak: E750-E752, E1091
\item En büyük kazancı kümeli ve gerçek (VLSI ve ulusal) örneklerde bekleyin; GNN'ye karşı fark
kısa bütçelerde en büyüktür.
\item Öğrenilmiş bir hedefleme skorunu bu vekile karşı, model süresi bütçeye dahil ve eşit
duvar-saatinde sınayın; sinyalin niteliğini ayrıca eşit vuruş sayısında ölçün. AUC gibi kestirim
ölçütleri arama faydasının yerine geçmez.
\item Arama içinde öğrenme kullanacaksanız krediyi kabul edilen iyileşmelere dayandırın
($\max(0,r)$ ya da kabul göstergesi), işaretli ortalama ödülden kaçının ve ağırlıkların iyi bir
statik ön-bilgiden ne kadar saptığını izleyin.
\item LK/LKH düzeyinde bir yerel arama kullanıyorsanız vuruş hedeflemeden kazanç beklemeyin; LKH-3'te
etkisi ölçülemedi.
\end{itemize}

\subsection{Kapsam ve geçerlilik sınırları}\label{sec:tartisma-sinir}

Bu çalışmanın doğrulayıcı bulguları, 2-opt/Or-opt tabanlı yüksek frekanslı ILS motoru, 10~sn
duvar-saati bütçesi ve $n=1.689$--$9.882$ ölçek aralığı için tesis edilmiştir.
Lin--Kernighan (LK) düzeyindeki yerel aramada hedeflemenin katkısı keşifsel olarak sınandı:
LKH-3'ün ILS kipinde hedefli vuruş rastgele vuruştan ayrışmadı. LKH'de yalnız tek bir yama, tek
bir bütçe ve ILS kipi denendi; bulgular 2-opt/Or-opt tabanlı hızlı ILS için geçerlidir.
Benzer biçimde segment vuruşu bu motorda sabit $W=50$ penceresiyle değerlendirilmiş olup, LK
tabanlı çözücülerdeki davranışı kendi mimari bağlamında ele alınmalıdır.
Keşifsel 60~sn bloğu (5 tohum), vekilin GNN'ye karşı avantajının daha uzun bütçelerde kademeli
olarak daralmakla birlikte korunduğunu göstermektedir ($-0{,}078$ puan, \%95 GA üst sınırı
$-0{,}008$); dakika mertebesini aşan bütçelerdeki asimptotik denge gelecek çalışmalarda
incelenebilir.
Donanım yükü dalgalanmalarının etkisini dışlamak amacıyla, bütün strateji kıyasları aynı oturum ve
aynı işçi sürecinde eşleşik olarak yürütülmüştür (ana oturumda rastgele global kol 3.252, yerel
vuruş oturumunda 2.225, ham uzunluk oturumunda 3.269 vuruş/sn; vekil sırasıyla 2.288 ve
2.773 vuruş/sn).
Kol sırası ana koşuda tohumlar arasında döndürüldü; ilk sıradaki kolun yaklaşık \%18 daha yavaş
vurduğu sonradan ölçüldü. Tam rastgele sırayla ve tek işçiyle koşulan bloklarda da vekil rastgeleden
düşük gap verdi.
% kaynak: E597, E598, E697, E705, E986, E987, E1057, E1078, E1113, E1164, E1181

Karşılaştırılan GNN modeli, $n\le400$ sentetik örneklerde eğitilmiş hafif bir mimari olarak büyük
test örneklerinde ($n\le9.882$) boyut genellemesi senaryosunu temsil eder; doğrudan hedef boyutta
eğitilmiş daha yüksek kapasiteli ağların hız--kalite takası ayrı bir değerlendirme konusudur.
Yayımlanmış ve daha büyük veriyle eğitilmiş NeuroLKH ile yapılan keşifsel karşılaştırma bu sınırı
kısmen kapatır: ağın süresi sayıldığında vekil ondan kötü değildir, sayılmadığında NeuroLKH
problemlerin çoğunda küçük bir farkla öndedir.
V4 testinde kazanç-MLP varsayılan havuz kuralıyla, vekil ise geliştirme setinde belirlenen kuralla
koşulmuştur; geliştirme setinde vekilin seçilen ve varsayılan kuralları arasında anlamlı bir fark
bulunmaması ($p=0{,}36$) bu karşılaştırmanın kararlılığını destekler. Kazanç-MLP aynı seçim
kuralıyla ayarlandığında da vekil önde kaldı ($-0{,}070$ puan).
Hedeflemenin etkinliği problem yapısına bağlıdır: kümeli ve gerçek problemlerde belirgin kazançlar
sağlanırken, kenar uzunluklarının homojen olduğu uniform dağılımlarda kazanç beklendiği üzere
sınırlı kalır ve tekil örneklerde (örn.\ h2c13) öğrenilmiş skorlar öne çıkabilir.
Son olarak, ön-kayıt protokolü (OpenTimestamps, Bitcoin blok 976258) uyarınca V1--V4 testleri
doğrulayıcı niteliktedir; katman analizleri, ikincil ikili karşılaştırmalar, ham uzunluk kolu ve
sağlamlık blokları ise bu ana bulguları derinleştiren ve sonraki ön-kayıtlı çalışmalar için
hipotez üreten keşifsel kanıtlar olarak konumlandırılmıştır.
% kaynak: E585, E751, E1107, E1110

\section{Sonuç}\label{sec:sonuc}

Hızlı bir ILS'te double-bridge vuruşlarının nereye yönlendirileceğini, öğrenilmiş ve eğitimsiz
sinyalleri ön-kayıtlı bir doğrulayıcı testte karşılaştırarak inceledik. Uygulayıcı için sonuç
basittir: vuruşları en uzun tur kenarlarına (ön-kayıtlı vekilde uç noktaların yerel ölçeğine göre
düzeltilmiş) yönlendirmek, bu motor ve 10~sn bütçede rastgele vuruştan geliştirme ve doğrulama
setlerinde tutarlı biçimde iyidir; sınanan öğrenilmiş skorlar karşısında bunun için model eğitmek
gerekmedi. Bu eğitimsiz vekil, LKH-3 turlarıyla eğitilmiş GNN'ye ve kazanç hedefli MLP'ye karşı
kötü-değil testlerini geçti, doğrulama setinde ortalama gap'i \%1,094'ten \%0,922'ye indirdi ve
keşifsel karşılaştırmalarda öğrenilmiş skorların önünde kaldı.
Bu başarım üç tamamlayıcı özellikle tutarlıdır: kazanca hizalı hedefleme, ihmal edilebilir
hesaplama maliyeti ve sürekli güncel havuz yapısı. Keşifsel bir kolda yerel ölçek düzeltmesi olmadan en
uzun kenarları hedeflemek de rastgeleye karşı benzer bir kazanç sağladı ve vekilden ortalamada
ayrışmadı ($+0{,}009$ puan); bu, yararın bu motorda doğrudan kenar uzunluğu bilgisinden geldiğini
düşündürür ve en basit vekile karşı sınama ilkesinin kendi önerimize de yararla uygulanabildiğini
gösterir.
Koşulardan önce kaydedilen keşifsel bloklar bu tabloyu genişletir: hedefli global vuruş, klasik
rastgele segment vuruşundan 1,04 puan iyidir; vekilin rastgeleye ve GNN'ye karşı avantajı 60~sn'de
daralarak sürer; geliştirmede seçilen havuz kuralının çevresindeki plato holdout2'de de görülür.
Hakem odaklı bloklar tabloyu netleştirir. Yayımlanmış NeuroLKH modeli, süresi bütçeden düşüldüğünde
vekilden ayrışmaz; süresi sayılmadığında problemlerin çoğunda küçük bir farkla öndedir. Vuruş sayısı
eşitlendiğinde öğrenilmiş skorlar vekilden kötü değildir ve vuruş başına en iyi sinyal eğitimsiz kNN
sırasıdır. Yani eşit süredeki sıralamayı sinyalin niteliğinden çok maliyeti belirler; vuruş başına en iyi
sinyal de, model süresi sayıldığında eşit sürede en düşük ortalama gap'i veren sinyal de eğitimsizdir. LKH-3'te ise hedefleme ölçülebilir bir kazanç getirmedi; bulgular hızlı,
2-opt/Or-opt tabanlı ILS için geçerlidir.
% kaynak: E608, E610, E965, E972, E1047, E1077, E1078, E1091, E1106, E1107, E1128-E1131, E1172

Koşu içi çevrimiçi öğrenme ise işaretli ödülle kendi ödülünü artırırken son kaliteyi düşürdü: vuruş
başına ortalama ödül, yerel aramanın onardığı nötr vuruşları ödüllendirir. Kredi kabul edilen
iyileşmelere dayandırıldığında bu zarar kayboldu, ama sabit modele bir kazanç da eklenmedi. Seyrek
iyileşmeli aramalarda kredi tasarımı bu nedenle öncelikli bir sorudur.
% kaynak: E1150-E1152 Yöntemsel ders de aynı yöndedir: öğrenilmiş bir arama bileşeni
rastgele tabana değil, en basit kazanca hizalı eğitimsiz vekile karşı ve eşit duvar-saatinde
değerlendirilmelidir.

Sonraki adımlar, karşılaştırmanın dakikalar mertebesini aşan bütçelerde tekrarı; vekil, kNN sırası
ve NeuroLKH arasındaki farkların ön-kayıtlı bir doğrulayıcı testle sınanması; ham ve ölçekle
düzeltilmiş uzunluğun ön-kayıtlı bir eşdeğerlik testiyle karşılaştırılması ve LK düzeyindeki
aramalarda hedeflemenin neden etkisiz kaldığının incelenmesidir.


\backmatter

\bmhead{Teşekkür}
Yok.

\section*{Beyanlar}

\bmhead{Fon desteği}
Yazar, bu çalışma için herhangi bir kurumdan fon desteği almamıştır.

\bmhead{Çıkar çatışması}
Yazarın bu makalenin içeriğiyle ilgili finansal ya da finansal olmayan bir çıkar çatışması
yoktur.

\bmhead{Etik onay}
Çalışmada insan ya da hayvan katılımcı ve kişisel veri kullanılmamıştır.

\bmhead{Yazar katkıları}
Çalışmanın kavramsallaştırılması, yazılım, deneyler, analiz ve yazım Ufuk ASIL tarafından
yürütülmüştür.

\bmhead{Yapay zekâ kullanımı}
Yazılım geliştirme süreçlerinde yardımcı kodlama araçları ve metnin dilsel düzenlemesinde büyük
dil modelleri destekleyici olarak kullanılmıştır. Deneysel tasarım, algoritmik doğrulama, tüm
istatistiksel analizler ve makalenin bilimsel içeriği bütünüyle yazarın denetiminde yürütülmüş
olup tüm sorumluluk yazara aittir.

\bmhead{Veri erişilebilirliği}
Waterloo VLSI, National ve TSPLIB örnekleri kamuya açıktır; sentetik setler ve holdout2, depodaki
betiklerle sabit tohumlardan yeniden üretilir. Dört ön-kayıt protokolü (Aşama~2, doğrulayıcı Aşama~3,
keşifsel sağlamlık blokları ve hakem odaklı keşifsel bloklar) SHA-256 özetleri, OpenTimestamps kanıtları ve sapma kaydıyla
birlikte; her (problem, tohum) koşusunun ham kaydı ve bütün analiz özetleri Zenodo'da kalıcı olarak
arşivlenmiştir (\href{https://doi.org/10.5281/zenodo.24396310}{doi:10.5281/zenodo.24396310}).
Doğrulayıcı protokolün zaman damgası Bitcoin blok 976258'e bağlıdır.
% kaynak: E584, E585

\bmhead{Kod erişilebilirliği}
Test ortamı, deney sürücüleri, model eğitimi, analiz betikleri ve tablo/şekil üreticileri MIT
lisansıyla \url{https://github.com/ufukasia/joco-kick-targeting} adresinde açıktır; kalıcı kopya
Zenodo'dadır (\href{https://doi.org/10.5281/zenodo.24396310}{doi:10.5281/zenodo.24396310}, sürüm
v1.2.0). Deponun README dosyası, her tablo ve şekli üreten komutları ve dosya düzenini listeler.
\path{scripts/verify_paper_numbers.py}, makaledeki başlıca sayıları ham koşu kayıtlarından yeniden
hesaplayarak doğrular.

\bibliography{references}

\end{document}
